\chapter{Apéndice experimental}
\label{sec:supplement}

Para cada capítulo se recogen aquí sus afirmaciones principales y aquello en lo que se sostienen: el experimento en el que se midió, qué se midió exactamente y dónde se publicó. El estado en la columna derecha indica cuán firme es la afirmación: \emph{medido}, hay un experimento directo; \emph{teoría}, una consecuencia matemática deducida en el capítulo o en un trabajo clásico; \emph{modelo}, el resultado de un cálculo con un modelo del libro (scripts en la carpeta \texttt{models/}); \emph{estimación}, un orden de magnitud sin medición directa; \emph{hipótesis}, una explicación plausible que el experimento aún no ha demostrado. Donde no hay experimento, se dice.

\section{Capítulo~\ref{sec:neurontypes}. Tipos de neuronas}

Las cifras del capítulo se apoyan en recuentos directos de células en el cerebro humano allí donde existen; la regla “una neurona, un neurotransmisor”, en atlas celulares y en experimentos que también encontraron excepciones; el papel de \nt{DOPA}, en registros de espigas; y los papeles de los demás canales de difusión, en hipótesis.

{\footnotesize
\setlength{\tabcolsep}{3pt}\renewcommand{\arraystretch}{1.15}
\begin{longtable}{>{\raggedright}p{0.5cm}>{\raggedright}p{4.0cm}>{\raggedright}p{5.6cm}>{\raggedright}p{2.3cm}>{\raggedright\arraybackslash}p{1.7cm}}
\toprule
N.º & Afirmación & Experimento y qué se midió & Fuente & Estado \\
\midrule\endfirsthead
\toprule
N.º & Afirmación & Experimento y qué se midió & Fuente & Estado \\
\midrule\endhead
1 & El cerebro humano tiene unas $8.6\cdot10^{10}$ neuronas, y casi todas las neuronas \nt{GLUT} son pequeñas neuronas del cerebelo (tabla~\ref{tab:types}) & Fraccionador isotrópico: el tejido se disuelve en una suspensión homogénea de núcleos y se cuentan los núcleos con el marcador neuronal NeuN. Los varones adultos tienen $86.1\pm8.1$ mil millones de neuronas; en la corteza cerebral está solo el $19\%$, y la mayor parte, en el cerebelo & \cite{Azevedo2009} & medido \\
2 & Casi toda neurona tiene un neurotransmisor principal (proposición~\ref{prop:dale}), pero hay excepciones & Atlas transcriptómico del cerebro del ratón: unos $4\cdot10^6$ células, $5322$ grupos; a cada grupo se le asigna un neurotransmisor por los genes de síntesis y empaquetado. Las excepciones se midieron directamente: las neuronas \nt{DOPA} en rebanadas liberan también GABA a través del mismo transportador vesicular e inhiben neuronas del estriado; en parte de los axones \nt{DOPA}, el glutamato y la dopamina salen de terminales distintas del mismo cable (microscopía electrónica y optogenética) & \cite{Strata1999,Yao2023,Tritsch2012,Zhang2015} & medido \\
3 & Toda la columna de la matriz de pesos tiene un solo signo~\eqref{eq:dalesign} & Deducción del capítulo a partir de la proposición~\ref{prop:dale} y de que casi todos los receptores de glutamato son excitadores y los de GABA, inhibidores. No hay una comprobación directa de “todos los destinatarios de una neurona reciben un peso del mismo signo” como regla general; contraejemplos: la coliberación de GABA por neuronas \nt{DOPA} (fila~2) y los destinatarios con receptores de signo distinto (fila~11) & deducción del capítulo & teoría \\
4 & Entre tres cuartas y cuatro quintas partes de las neuronas de la corteza son \nt{GLUT}, entre una quinta y una cuarta parte, \nt{GABA} & Inmunotinción de GABA en columnas de $50$\,\textmu m de ancho a través de todo el espesor cortical en $10$ áreas de la corteza de mono: en $8$ áreas las neuronas GABA son cerca del $25\%$ de todas las neuronas; en la corteza visual primaria, menos del $20\%$ & \cite{Hendry1987} & medido \\
5 & Una neurona \nt{GLUT} tiene unas $10^4$ salidas & Estereología post mortem: en la neocorteza de varones jóvenes hay $1.64\cdot10^{14}$ sinapsis (microscopía electrónica, disector) y unas $2.3\cdot10^{10}$ neuronas. El cociente da $\sim7\cdot10^3$ sinapsis por neurona; en promedio, el número de entradas es igual al de salidas & \cite{Tang2001synapses,Pakkenberg1997} & medido \\
6 & La salida de una neurona \nt{DOPA} se ramifica en $10^5$--$10^6$ terminales; es una estimación por cobertura de la diana, no un recuento & Marcaje viral de la membrana de neuronas \nt{DOPA} individuales de rata y reconstrucción completa del axón: una neurona cubre entre el $0.45$ y el $5.7\%$ (en promedio el $2.7\%$) del volumen del estriado. Se midió la cobertura, no el número de terminales: este se dedujo de la cobertura en orden de magnitud, sin recuento directo. Para \nt{SERO} y \nt{NORA}, los números de terminales son estimaciones burdas & \cite{Matsuda2009} & medido \\
7 & Las neuronas de difusión son pocas: \nt{DOPA} $\sim5\cdot10^5$ (el principal de dos grupos, cota inferior), \nt{SERO} $\sim3\cdot10^5$ (suma de dos recuentos parciales), \nt{HIST} $\sim6\cdot10^4$ (tabla~\ref{tab:types}, fig.~\ref{fig:typepie}) & Recuentos estereológicos en humanos: $550\,000$ neuronas pigmentadas en la sustancia negra (sin el área tegmental ventral); $235\,000$ neuronas en el núcleo dorsal del rafe (todas las del núcleo) y otras $125\,000$ neuronas serotoninérgicas en el tegmento pontino; unas $64\,000$ neuronas histaminérgicas del hipotálamo. Para \nt{NORA}, \nt{GLYC}, \nt{ACET} y \nt{ADRE} no hay recuento directo: en la tabla son estimaciones burdas del orden de magnitud & \cite{Pakkenberg1991,Baker1990,Baker1991,Airaksinen1991} & medido \\
8 & El glutamato en la hendidura sube hasta cerca de un milimolar y decae en $\sim1\ms$ & Por el desplazamiento de un bloqueador competitivo de disociación rápida de los receptores NMDA durante la transmisión sináptica (cultivo de neuronas de hipocampo): pico de $1.1$\,mM, constante de decaimiento de $1.2\ms$ & \cite{Clements1992} & medido \\
9 & En la corteza activa, las corrientes excitadora e inhibidora están casi equilibradas y la neurona se mantiene cerca del umbral & Teoría: una red con conexiones fuertes llega por sí sola a un régimen equilibrado. Experimento: en la corteza prefrontal del hurón in vivo, durante los estados “up”, el potencial de inversión de la corriente sináptica se mantiene cerca de $-37$\,mV durante cientos de milisegundos, aunque la conductancia cambia en $\sim21$\,nS: la excitación y la inhibición suben y bajan juntas & \cite{vanVreeswijk1996,Haider2006} & medido \\
10 & Las neuronas \nt{DOPA} comunican el error de predicción de la recompensa~\eqref{eq:rpe} (fig.~\ref{fig:rpe}) & Espigas de neuronas \nt{DOPA} de mono: ráfaga ante el jugo inesperado; tras el aprendizaje, ráfaga ante la señal y ninguna respuesta al jugo predicho; caída de la frecuencia cuando no se da el jugo prometido. En el experimento cuantitativo, la frecuencia es proporcional a la diferencia entre la recompensa actual y una media ponderada de las pasadas, pero solo para errores positivos. La ecuación~\eqref{eq:rpe} misma es el algoritmo de diferencias temporales & \cite{Schultz1997,Bayer2005,SuttonBarto2018} & medido \\
11 & El signo y la velocidad los fija el receptor: los ionotrópicos, fracciones de milisegundo; los metabotrópicos, mucho más lentos (proposición~\ref{prop:receptor}) & Pulsos rápidos de glutamato sobre parches de membrana de neuronas de hipocampo: la corriente por los receptores ionotrópicos sube (del $20$ al $80\%$) en $0.2$--$0.6\ms$ y decae en $2$--$3\ms$. El potencial inhibidor lento de las neuronas piramidales se suprime con un bloqueador selectivo del receptor metabotrópico de GABA. Dos familias de receptores de dopamina con acción opuesta; en el estriado, las neuronas con receptor D1 necesitan dopamina para reforzar sinapsis, y las neuronas con D2, para debilitarlas & \cite{Colquhoun1992,DutarNicoll1988,Beaulieu2011,Shen2008} & medido \\
12 & Un canal de difusión no es un solo cable, sino un mazo de pocas líneas (sección~\ref{sec:buses}) & Marcaje genético-viral de las neuronas serotoninérgicas del núcleo dorsal del rafe del ratón: las neuronas que envían su salida a la amígdala y a la corteza frontal ocupan lugares distintos del núcleo, se ramifican en regiones complementarias y responden de forma opuesta a un estímulo desagradable. Revisión: subgrupos de neuronas del locus coeruleus (\nt{NORA}) codifican procesos distintos & \cite{Ren2018,Poe2020} & medido \\
13 & \nt{NORA} fija la ganancia $g$ & Hipótesis de la ganancia adaptativa; modelo de red en el que las catecolaminas aumentan la pendiente de la característica de transferencia. No hay una medición directa de “$g$ crece con la noradrenalina” en toda la red; se midió que el diámetro de la pupila sigue las espigas de las neuronas del locus coeruleus del mono & \cite{AstonJones2005,ServanSchreiber1990,Joshi2016} & hipótesis \\
14 & \nt{ACET} conmuta “escritura/lectura” y fija la tasa de aprendizaje & Hipótesis. Apoyo: en rebanadas de corteza olfativa de rata, los agonistas de la acetilcolina ($100$\,\textmu M) debilitan las sinapsis internas, las que “recuerdan”, en más del $60\%$, y las de entrada, en menos del $15\%$ & \cite{Hasselmo2006,HasselmoBower1992} & hipótesis \\
15 & \nt{SERO} fija el factor de descuento $\gamma$; las neuronas serotoninérgicas trabajan de forma regular, unas pocas espigas por segundo, y casi callan en una de las fases del sueño & Hipótesis “serotonina $=\gamma$”. El argumento más fuerte: la activación optogenética de las neuronas serotoninérgicas del núcleo dorsal del rafe del ratón reduce el número de abandonos prematuros de la espera de una recompensa diferida. La frecuencia de espigas y el silencio durante el sueño de movimientos oculares rápidos están medidos (revisión de registros) & \cite{Doya2002,Miyazaki2014,JacobsAzmitia1992} & hipótesis \\
\bottomrule
\end{longtable}}

\section{Capítulo~\ref{sec:glutcomputer}. Qué calculan las neuronas \nt{GLUT}}

Las proposiciones lógicas del capítulo son teoremas sobre circuitos de pesos no negativos, deducidos en el propio capítulo; los experimentos confirman el modelo de neurona y que los circuitos construidos según esos teoremas (detector de coincidencias, línea de retardo, grupo autosostenido, cadena) existen realmente en el cerebro.

{\footnotesize
\setlength{\tabcolsep}{3pt}\renewcommand{\arraystretch}{1.15}
\begin{longtable}{>{\raggedright}p{0.5cm}>{\raggedright}p{4.0cm}>{\raggedright}p{5.6cm}>{\raggedright}p{2.3cm}>{\raggedright\arraybackslash}p{1.7cm}}
\toprule
N.º & Afirmación & Experimento y qué se midió & Fuente & Estado \\
\midrule\endfirsthead
\toprule
N.º & Afirmación & Experimento y qué se midió & Fuente & Estado \\
\midrule\endhead
1 & La neurona es un circuito RC con umbral y reinicio~\eqref{eq:glutneuron} & Al soma de neuronas piramidales de corteza de rata (rebanadas) se aplicó una corriente ruidosa parecida al flujo sináptico del cerebro vivo. La dependencia de la frecuencia media respecto de la media y la dispersión de la corriente se describe con la neurona integradora con fugas si se añade una adaptación lenta de la frecuencia. Los intervalos de $\tau$ y de los retardos se dan en el capítulo sin referencia experimental & \cite{Rauch2003} & medido \\
2 & El modelo~\eqref{eq:glutneuron} es una simplificación: las ramas de entrada generan por sí mismas espigas locales & Registro directo en las dendritas de neuronas piramidales de la corteza visual del ratón in vivo: un estímulo visual provoca espigas dendríticas locales dependientes de los receptores NMDA; su supresión ensancha la sintonía de la neurona a la orientación & \cite{Smith2013} & medido \\
3 & La ventana de coincidencia es $\Delta_{\max}=\tau\ln\frac{w}{\theta-w}$~\eqref{eq:window}; con $\theta=1.5w$ vale $0.7\tau$ & Consecuencia del modelo~\eqref{eq:glutneuron}. Hay una medición directa de una ventana de suma estrecha en neuronas con inhibición rápida (capítulo~\ref{sec:gaba}, fila~3 de su tabla) & deducción del capítulo & teoría \\
4 & Una red formada solo por neuronas \nt{GLUT} calcula únicamente funciones monótonas de las entradas (proposición~\ref{prop:monotone}) & Demostración en el capítulo: iteración desde el silencio con un lado derecho no decreciente. Es una proposición sobre el modelo; no hay un experimento que la haya comprobado en una red viva sin inhibición & deducción del capítulo & teoría \\
5 & Los órganos de los sentidos entregan un par de cables: unas células responden al aumento del estímulo y otras a su disminución & Retina: ya en las células bipolares la señal de los conos se divide en canales de encendido y de apagado. Bloqueo farmacológico de las bipolares de encendido en el mono: los canales van separados hasta la corteza; tras el bloqueo empeora la detección del aumento de luz, pero no la de su disminución & \cite{Schiller1992} & medido \\
6 & Con código de dos cables, una red de neuronas \nt{GLUT} calcula cualquier función lógica de forma asíncrona, sin ciclo de reloj común, al precio del doble de neuronas (proposición~\ref{prop:dualrail}, \eqref{eq:dualrail}, fig.~\ref{fig:dualxor}) & El código de dos cables y la detección de la finalización por la propia señal son una técnica estándar de los circuitos asíncronos insensibles a los retardos. El circuito de OR exclusivo de seis neuronas se construye en el capítulo. No hay experimentos que muestren que el cerebro calcula funciones lógicas precisamente con código de dos cables & \cite{Sparso2001} & teoría \\
7 & La mayor diferencia de tiempo de llegada del sonido a los oídos en el ser humano es $\approx0.6\ms$, y el cerebro distingue unos diez microsegundos & Oyentes en un experimento de elección forzada entre dos alternativas: el umbral de discriminación de la diferencia de tiempo (75\% de aciertos) es de $9$\,\textmu s para ruido en la banda de $150$--$1700$\,Hz, $11$\,\textmu s para un tono de $1$\,kHz y $28$\,\textmu s para un clic. El límite de $0.6\ms$ es un cálculo del capítulo, $0.22\,\text{m}/343\,\text{m/s}$ & \cite{KlumppEady1956} & medido \\
8 & En las aves, la diferencia de tiempos se convierte en lugar mediante líneas de retardo y detectores de coincidencias~\eqref{eq:itd} (fig.~\ref{fig:itd}) & Lechuza común: los axones de los dos oídos entran en el núcleo de detectores de coincidencias por lados opuestos; el tiempo de llegada de las espigas a lo largo del axón cambia de forma ordenada con la profundidad, y el rango de retardos a través del núcleo ($160$\,\textmu s en $700$\,\textmu m) coincide con el rango de retardos interaurales & \cite{Carr1990} & medido \\
9 & En los mamíferos no se ha encontrado una fila de retardos; la misma tarea la resuelven detectores de coincidencias con una inhibición de \nt{GLYC} temporizada con precisión & Registro in vivo en la oliva superior medial del jerbo: las respuestas no forman un mapa de direcciones; la aplicación de glicina y de su bloqueador por micropipeta muestra que una inhibición \emph{glicinérgica} temporizada con precisión fija el rango de trabajo de los retardos. Aquí la inhibición proviene de \nt{GLYC}, no de \nt{GABA} & \cite{Brand2002,Grothe2010} & medido \\
10 & Un grupo con realimentación fuerte es un biestable; la actividad autosostenida dura segundos tras el estímulo (fig.~\ref{fig:bistable}) & Corteza prefrontal de mono en una tarea de movimiento ocular diferido hacia un punto memorizado (retardo de $1$--$6$\,s): $87$ de $170$ neuronas relacionadas con la tarea cambian su frecuencia durante el retardo, y en el $79\%$ de ellas esto depende de la dirección del punto memorizado. Que esta actividad la sostenga una realimentación con dos estados estables es teoría & \cite{Funahashi1989,Wang2001} & medido \\
11 & El bit se escribe si el aporte dura más de $\approx0.7\tau$ (fig.~\ref{fig:recurrentgroup}b) & Cálculo de la ecuación $\tau\,\dd r/\dd t=-r+f(wr+I)$ por el método de Euler con $w=1$, $I=0.6$; no hay script en \texttt{models/}. No hay experimento & cálculo del capítulo & modelo \\
12 & La memoria puede guardarse en la facilitación sináptica, casi sin espigas & Teoría: el calcio residual en las terminales mantiene lo escrito cerca de un segundo. Apoyo: en la corteza prefrontal medial del hurón (registro de varias neuronas a la vez) hay una subred de neuronas piramidales conectadas por sinapsis facilitadoras con un efecto residual prolongado. Revisión de observaciones de intervalos “silenciosos” durante el mantenimiento en memoria. Qué forma usa la corteza y cuándo es una pregunta abierta & \cite{Mongillo2008,WangMarkram2006,Stokes2015} & hipótesis \\
13 & La memoria se borra por fatiga: se agota la reserva de neurotransmisor & Registros por pares de neuronas piramidales de la corteza: la respuesta de la sinapsis cae con espigas frecuentes, y la velocidad de la caída la fija la probabilidad de liberación. Que sea precisamente esto lo que apaga un grupo autosostenido no se ha mostrado directamente & \cite{Tsodyks1997} & medido \\
14 & Una cadena de grupos es un temporizador disparado por un evento; en las aves canoras las neuronas se disparan estrictamente por turno & Registro de neuronas identificadas del centro vocal superior del diamante mandarín durante el sueño y el canto: cada neurona que envía su salida al núcleo motor emite una sola ráfaga breve en un instante preciso del canto, y las neuronas se disparan en secuencia & \cite{Hahnloser2002} & medido \\
\bottomrule
\end{longtable}}

\section{Capítulo~\ref{sec:gaba}. \nt{GLUT} y \nt{GABA} juntos}

Las proposiciones lógicas y dinámicas del capítulo se deducen en él mismo a partir del análisis lineal y de la ley de Ohm; los experimentos muestran que los circuitos en los que se apoyan (el veto con retardo, la inhibición directa tras la excitación, la estabilización por inhibición, el ritmo del lazo \nt{GLUT}--\nt{GABA}) funcionan realmente en el cerebro.

{\footnotesize
\setlength{\tabcolsep}{3pt}\renewcommand{\arraystretch}{1.15}
\begin{longtable}{>{\raggedright}p{0.5cm}>{\raggedright}p{4.0cm}>{\raggedright}p{5.6cm}>{\raggedright}p{2.3cm}>{\raggedright\arraybackslash}p{1.7cm}}
\toprule
N.º & Afirmación & Experimento y qué se midió & Fuente & Estado \\
\midrule\endfirsthead
\toprule
N.º & Afirmación & Experimento y qué se midió & Fuente & Estado \\
\midrule\endhead
1 & Las neuronas \nt{GABA} son entre una quinta y una cuarta parte de las de la corteza & Inmunotinción de GABA en $10$ áreas de la corteza de mono: cerca del $25\%$ de las neuronas en $8$ áreas, menos del $20\%$ en la corteza visual primaria. Revisión de la diversidad de las neuronas inhibidoras de la corteza & \cite{Hendry1987,Markram2004} & medido \\
2 & Lo que hace la inhibición depende en gran medida del lugar de llegada: dendrita, cuerpo, lugar de nacimiento de la espiga, terminal de otro cable & Revisión: las clases de neuronas \nt{GABA} de la corteza se distinguen por su diana en el destinatario. Registro simultáneo del cuerpo y de una dendrita de una neurona piramidal: la inhibición directa es mucho más fuerte en el cuerpo que en las dendritas. La inhibición sobre la terminal de otro cable en la médula espinal reduce la liberación de neurotransmisor de una sola línea de entrada (revisión de mediciones) & \cite{Markram2004,Pouille2001,Rudomin1999} & medido \\
3 & El cambio de signo a través de un intermediario cuesta un retardo, cerca de un milisegundo; la inhibición que llega tras la excitación estrecha la ventana de coincidencia & Neuronas piramidales de hipocampo de rata en rebanadas: la inhibición a través de un intermediario llega tras la excitación directa con un retardo corto, y la ventana en la que dos entradas excitadoras se suman hasta el umbral es menor de $2\ms$. En las dendritas, donde esta inhibición es más débil, la ventana es más ancha & \cite{Pouille2001} & medido \\
4 & El veto $a\wedge\bar b$~\eqref{eq:veto} con OR y una constante uno es funcionalmente completo (proposición~\ref{prop:gabacomplete}) & Demostración en el capítulo. Resultado clásico: una red de elementos de umbral con entradas inhibidoras realiza cualquier expresión lógica. Proposición sobre el modelo; no hay experimento & \cite{McCullochPitts1943} & teoría \\
5 & El veto con retardo da un detector de dirección del movimiento con velocidad óptima $v^*=\Delta x/d$~\eqref{eq:vstar} & Retina de conejo: la selectividad a la dirección la crea una inhibición que, en el movimiento en la dirección “nula”, apaga la excitación. El registro por separado de las entradas excitadora e inhibidora de las células selectivas mostró que las células amacrinas estrelladas (\nt{GABA}) las inhiben de forma asimétrica, solo con las prolongaciones dirigidas hacia el lado “nulo”. La fórmula de $v^*$ es una deducción del capítulo & \cite{Barlow1965,Fried2002} & medido \\
6 & La inhibición de derivación divide el voltaje, no le resta~\eqref{eq:shunt} (fig.~\ref{fig:shunt}) & Ley de Ohm para un divisor de tres conductancias con el potencial de equilibrio del cloro cerca del reposo; para una neurona sin espigas & deducción del capítulo & teoría \\
7 & Sobre la frecuencia de espigas, la derivación actúa restando, y la división surge de la inhibición junto con un ruido de fondo equilibrado & Modelo: en una neurona que se dispara, la corriente por la derivación es casi constante y el efecto sobre la frecuencia es de resta. Experimento: en neuronas piramidales de corteza somatosensorial de rata (rebanadas) se inyectaron mediante pinza dinámica conductancias de fondo excitadoras e inhibidoras; el aumento simultáneo y equilibrado de ambas divide la pendiente de la respuesta sin un desplazamiento apreciable de la frecuencia. Revisión: la división por la actividad total se encuentra en muchas regiones del cerebro & \cite{HoltKoch1997,Chance2002,Carandini2012} & medido \\
8 & Con $\beta>1$, la inhibición mutua elige un solo ganador (proposición~\ref{prop:wta}, \eqref{eq:wta}) & Los valores propios $-(1\pm\beta)/\tau$ son una deducción del capítulo. Las frecuencias $0.73$ y $0.53$ con $\beta=0.5$ son el punto fijo $r_{1,2}=(I_{1,2}-\beta I_{2,1})/(1-\beta^2)$; no hay script en \texttt{models/}. El capítulo no cita un experimento directo & deducción del capítulo & teoría \\
9 & Reinicio de la memoria con una señal a través de \nt{GABA}; dos grupos con inhibición mutua forman un cerrojo & Se sigue de la fig.~\ref{fig:bistable} y de la proposición~\ref{prop:wta}. No hay experimento & deducción del capítulo & teoría \\
10 & El equilibrio del lazo \nt{GLUT}--\nt{GABA} es estable si la inhibición es suficientemente fuerte y suficientemente rápida (proposición~\ref{prop:ei}) & Modelo de dos poblaciones~\eqref{eq:ei}; las condiciones (i) y (ii) son el determinante y la traza de su matriz, deducidas en el capítulo & \cite{WilsonCowan1972} & teoría \\
11 & Con $w_{EE}=1.5$ y $w_{EI}w_{IE}=2$, la inhibición rápida ($\tau_I=2\ms$) mantiene $E^*=0.67h$, y la lenta ($30\ms$) desencadena oscilaciones (fig.~\ref{fig:ei}) & Solución numérica de~\eqref{eq:ei}, script \texttt{figures/make\_ei.py} & cálculo & modelo \\
12 & La corteza funciona en régimen de estabilización por inhibición; su firma: si se empuja a las neuronas \nt{GABA}, su frecuencia baja~\eqref{eq:isn} & La teoría predijo la respuesta paradójica. Experimento: registros intracelulares en la corteza visual primaria del gato; al suprimir la respuesta con un estímulo circundante, la conductancia inhibidora no crece, sino que cae junto con la excitadora. La excitación optogenética de neuronas PV del ratón en las cortezas visual, somatosensorial y motora reduce la frecuencia de las neuronas inhibidoras, también sin estímulo y con anestesia ligera. El coeficiente $-1/3$ con los números de la fig.~\ref{fig:ei} es una deducción del capítulo & \cite{Tsodyks1997isn,Ozeki2009,Sanzeni2020} & medido \\
13 & El lazo “\nt{GLUT} excita a \nt{GABA}, \nt{GABA} inhibe a \nt{GLUT}” genera un ritmo rápido con período de $12$--$50\ms$ & La estimulación optogenética de neuronas \nt{GABA} rápidas en la corteza somatosensorial del ratón a frecuencias de $8$--$200$\,Hz refuerza selectivamente el ritmo gamma ($20$--$80$\,Hz, período de $12$--$50\ms$); la estimulación de neuronas \nt{GLUT} refuerza solo ritmos más lentos & \cite{Cardin2009,Buzsaki2012} & medido \\
\bottomrule
\end{longtable}}

\section{Capítulo~\ref{sec:code}. Código Morse}

Las estimaciones límite del capítulo son teoría de la información y cálculo; está medido dónde surge el ruido en el canal, qué tan rápido trabaja el cerebro y cuánto cuesta una espiga, mientras que la pregunta de qué código usa realmente la corteza sigue abierta.

{\footnotesize
\setlength{\tabcolsep}{3pt}\renewcommand{\arraystretch}{1.15}
\begin{longtable}{>{\raggedright}p{0.5cm}>{\raggedright}p{4.0cm}>{\raggedright}p{5.6cm}>{\raggedright}p{2.3cm}>{\raggedright\arraybackslash}p{1.7cm}}
\toprule
N.º & Proposición & Experimento y qué se midió & Fuente & Estatus \\
\midrule\endfirsthead
\toprule
N.º & Proposición & Experimento y qué se midió & Fuente & Estatus \\
\midrule\endhead
1 & Las frecuencias de las neuronas \nt{GLUT} de la corteza van de fracciones a decenas de espigas por segundo, y la mayor parte del tiempo las neuronas trabajan cerca del límite inferior & Registro con una pipeta adosada a la célula (la elección de la neurona no depende de su actividad) en la corteza auditiva de la rata despierta: en cada momento, menos del $5\%$ de las neuronas supera las $20$ espigas por segundo; en parte de las neuronas la frecuencia de fondo es menor de $0.01$ por segundo; la distribución de frecuencias es lognormal & \cite{Hromadka2008} & medido \\
2 & La capacidad del canal de espigas es $R_{\max}\approx r\log_2\frac{e}{r\Delta t}$~\eqref{eq:spikeentropy}, y casi toda está en el tiempo de las espigas (proposición~\ref{prop:capacity}) & Entropía de una cadena binaria de celdas de longitud $\Delta t$. Cifras del capítulo: $8$ bits por espiga con $r=10$ por segundo y $\Delta t=1\ms$, unos $5$ bits con $\Delta t=10\ms$, menos de $2$ bits por segundo para un contador de espigas & \cite{MacKay1952} & teoría \\
3 & Las neuronas sensoriales transmiten de uno a varios bits por espiga, en los mejores casos hasta la mitad del límite & Reconstrucción del estímulo a partir de las espigas de neuronas sensoriales cuyo estímulo se conoce; resumen de mediciones en el libro & \cite{Rieke1997} & medido \\
4 & El generador de espigas es preciso; el tiempo exacto lo fijan los cambios rápidos de la entrada & Neuronas de corteza de rata en rebanadas: con una corriente repetida de fluctuaciones rápidas, las espigas se repiten con precisión mejor que $1\ms$; con corriente constante, los instantes se dispersan & \cite{Mainen1995} & medido \\
5 & La sinapsis es poco fiable: más de la mitad de las espigas no llega al receptor por el azar de la liberación & Estimulación mínima de las entradas de neuronas piramidales del hipocampo (rebanadas): más de la mitad de las espigas no provoca respuesta. Se descartaron las fluctuaciones del umbral de estimulación, los fallos de conducción en las ramificaciones del axón y la baja temperatura del experimento; queda la liberación probabilística & \cite{AllenStevens1994} & medido \\
6 & El flujo de espigas en la corteza viva parece aleatorio: los intervalos se dispersan como en un flujo de Poisson y la varianza del número de espigas es cercana a la media; parte del “ruido” es un mensaje sobre los movimientos del animal & Registros en las áreas visuales V1 y MT del mono despierto: coeficiente de variación de los intervalos cercano a $1$, razón varianza/media del número de espigas como en un proceso aleatorio; un modelo que suma entradas pequeñas independientes da una dispersión mucho menor. En la corteza visual del ratón, una parte considerable de la dispersión se predice a partir de un video de los movimientos del propio animal & \cite{Softky1993,Stringer2019} & medido \\
7 & El código de frecuencia es lento: error $1/\sqrt{rT}$~\eqref{eq:rateerror}, mientras que el cerebro reconoce una imagen en $\sim150\ms$ & La fórmula es estadística de Poisson, deducción del capítulo. Experimento: personas decidían si había un animal en una fotografía desconocida mostrada durante $20\ms$; los potenciales cerebrales para “sí” y “no” divergen al cabo de unos $150\ms$. La división de ese tiempo en una decena de etapas de $10$--$20\ms$ es una estimación del capítulo, no una medición & \cite{Thorpe1996} & medido \\
8 & El promediado sobre un grupo está limitado por la correlación: el error cae como mucho $1/\sqrt c$ veces~\eqref{eq:correlated} (proposición~\ref{prop:pooling}) & Pares de neuronas del área MT del mono registrados a la vez: coeficiente de correlación de los números de espigas de $0.12$ en promedio, y el análisis teórico muestra que incluso esa correlación limita la ganancia del grupo. Teoría: el límite lo pone solo la parte de la correlación que se parece a la señal. Modelo de la postura contraria: la frecuencia de un grupo de $50$--$100$ neuronas se lee en un solo intervalo entre espigas, $10$--$50\ms$, y el tiempo de las espigas individuales no se usa en la corteza & \cite{Zohary1994,MorenoBote2014,ShadlenNewsome1998} & medido \\
9 & Un número puede escribirse en el tiempo de la primera espiga~\eqref{eq:latency}, en el orden de las espigas de un grupo o en la fase de un ritmo local & Retina de salamandra: la estructura espacial de una imagen mostrada brevemente está escrita en las latencias relativas de las primeras espigas de las neuronas de salida, casi independientemente del contraste. Células de lugar del hipocampo de rata: al atravesar “su” lugar, las espigas se adelantan cada vez más en el ciclo del ritmo theta ($7$--$12$\,Hz), con un desplazamiento de $100^\circ$ a $355^\circ$. Hasta qué punto usa la corteza el tiempo de las espigas es una pregunta abierta (fila~8) & \cite{Gollisch2008,OKeefe1993} & medido \\
10 & Qué código se lee lo decide la constante de tiempo del receptor~\eqref{eq:reader} (proposición~\ref{prop:reader}); en la corteza activa, las neuronas están más cerca de los detectores de coincidencias & La fórmula~\eqref{eq:reader} es una deducción del capítulo. Revisión de registros in vivo: en la corteza activa la conductancia de la membrana es varias veces mayor que en reposo, y la constante de tiempo, correspondientemente más corta. Que la corteza cambie de código precisamente así no se ha mostrado directamente & \cite{Destexhe2003} & hipótesis \\
11 & La sinapsis con depresión transmite el cambio de frecuencia, en esencia $\Delta r/r$~\eqref{eq:depression} (fig.~\ref{fig:depression}) & Registros por pares de neuronas piramidales de la corteza: la velocidad de la depresión la fija la probabilidad de liberación; con depresión lenta la respuesta refleja la frecuencia, y con depresión rápida, las coincidencias. Modelo basado en la depresión medida en la corteza de rata: cambios relativos iguales de frecuencia en entradas rápidas y lentas dan la misma respuesta, es decir, la sinapsis transmite $\Delta r/r$. Las cifras $1.7\to2.2$, $6.7$ y $90\ms$ son un cálculo del capítulo & \cite{Tsodyks1997,Abbott1997depression} & medido \\
12 & La ráfaga es una “raya”: su probabilidad de perderse, $(1-p)^k$~\eqref{eq:burst}, es pequeña, y la facilitación la reduce aún más & Revisión: en sinapsis poco fiables, las espigas aisladas se pierden a menudo, mientras que las ráfagas se transmiten de forma fiable gracias a la facilitación; una ráfaga provoca cambios sinápticos duraderos. Que el cerebro lea la ráfaga como un signo aparte es una hipótesis & \cite{Lisman1997,AllenStevens1994} & hipótesis \\
13 & Las neuronas \nt{GABA} rápidas fijan el ritmo y la “puntuación”; otra clase inhibe a las inhibidoras y abre la compuerta (proposición~\ref{prop:punctuation}) & Revisión de las propiedades de las clases de neuronas \nt{GABA} de la corteza. Optogenética: la excitación rítmica de las neuronas \nt{GABA} rápidas provoca el ritmo gamma, y la respuesta a un estímulo depende de la fase del ciclo. En la corteza visual del ratón, las neuronas PV se inhiben entre sí; las SST inhiben a todas las demás clases; las VIP, sobre todo a las SST. La excitación de las neuronas VIP en el ratón despierto retira la inhibición de las neuronas \nt{GLUT}; las activa una señal de refuerzo. La división de papeles como regla general es una hipótesis & \cite{Tremblay2016,Cardin2009,Pfeffer2013,Pi2013} & hipótesis \\
14 & La energía limita la frecuencia media a un valor del orden de una espiga por segundo~\eqref{eq:energy} & Presupuesto energético de la sustancia gris de roedor a partir de datos anatómicos y fisiológicos: las espigas y las corrientes sinápticas suponen el $47\%$ y el $34\%$ de la energía de señalización; no pueden estar activas a la vez más de $\sim15\%$ de las neuronas. Para la corteza humana: no más de $\sim1\%$ de las neuronas pueden estar muy activas a la vez. La cifra de $\sim10^9$ ATP por espiga y la potencia de $\sim10$\,W para la señalización son una estimación del capítulo basada en estos cálculos & \cite{Attwell2001,Lennie2003} & teoría \\
15 & El código es disperso y está ajustado al máximo de bits por julio; la corteza cambia precisión por energía (proposición~\ref{prop:economy}) & Hipótesis del código económico. Apoyo: en ratones con restricción de alimento, en la corteza visual caen la conductancia de los receptores AMPA y el gasto sináptico de ATP ($-29\%$), la frecuencia de espigas se mantiene y la sintonía a la orientación se ensancha un $32\%$ & \cite{Barlow1961,Padamsey2022} & hipótesis \\
\bottomrule
\end{longtable}}

\section{Capítulo~\ref{sec:serocomputer}. Neuronas \nt{SERO}}

Las propiedades de las propias neuronas \nt{SERO} (ritmo, autoinhibición, signo según el receptor, subsistemas) están medidas directamente; los cálculos del capítulo (regulador, filtro, lógica) se siguen de sus ecuaciones; $\tau_c$, el coeficiente de difusión de la serotonina y el número de sitios de liberación son estimaciones; el papel del lazo como regulador de nivel en el cerebro es una hipótesis (en el núcleo del rafe la autoinhibición es puntual y solo da pausas breves), y el papel de \nt{SERO} como horizonte es una hipótesis con experimentos de conducta a su favor.

{\footnotesize
\setlength{\tabcolsep}{3pt}\renewcommand{\arraystretch}{1.15}
\begin{longtable}{>{\raggedright}p{0.5cm}>{\raggedright}p{4.0cm}>{\raggedright}p{5.6cm}>{\raggedright}p{2.3cm}>{\raggedright\arraybackslash}p{1.7cm}}
\toprule
N.º & Proposición & Experimento y qué se midió & Fuente & Estado \\
\midrule\endfirsthead
\toprule
N.º & Proposición & Experimento y qué se midió & Fuente & Estado \\
\midrule\endhead
1 & El signo de la acción de la serotonina lo elige el receptor del destinatario (proposición~\ref{prop:receptor}) & Los receptores de serotonina se dividen en familias según su farmacología y su mecanismo: 5-HT$_{1A}$ abre canales de potasio a través de una proteína G e inhibe la célula; 5-HT$_{2A}$ y el ionotrópico 5-HT$_3$ excitan. Un mismo neurotransmisor da respuestas opuestas en células distintas. & \cite{BarnesSharp1999} & medido \\
2 & La neurona \nt{SERO} en vigilia emite de forma regular varias espigas por segundo & Registro de neuronas del núcleo dorsal del rafe en gatos en libertad de movimiento: descarga rítmica lenta de $2.8\pm0.2$ espigas/s en vigilia tranquila; en el sueño lento la frecuencia cae un $34$--$68\%$, en el sueño REM, un $98\%$. En ratas anestesiadas, $1.7\pm0.5$ espigas/s, muy regular. & \cite{TrulsonJacobs1979, GagnonParent2014, JacobsAzmitia1992} & medido \\
3 & La neurona \nt{SERO} es un marcapasos: no necesita un empujón para cada espiga, pero sí un aporte de fondo constante (en la rebanada, noradrenalina a través de receptores $\alpha_1$; en el organismo, desde \nt{NORA}) & En la rebanada de cerebro las células descargan de forma lenta y regular; tras cada espiga hay una poshiperpolarización de $200$--$800\ms$. En la rebanada hay menos células activas espontáneamente que en el organismo: el ritmo regular necesita una entrada tónica de noradrenalina a través de receptores $\alpha_1$, cuyo bloqueo lo apaga. & \cite{VandermaelenAghajanian1983} & medido \\
4 & Casi todos los receptores son metabotrópicos; la respuesta es lenta: sube y baja en cerca de un segundo & Todas las familias de receptores salvo 5-HT$_3$ están acopladas a proteínas G. En la rebanada, la liberación evocada de serotonina produce a través de 5-HT$_{1A}$ una corriente inhibidora que sube y baja en un segundo. & \cite{BarnesSharp1999, CourtneyFord2016} & medido \\
5 & En las regiones de destino la serotonina sale al volumen, no solo a la hendidura; en el propio núcleo del rafe la inhibición de las vecinas es punto a punto & Voltamperometría cíclica rápida en rebanadas: la liberación por espiga es igual en una región con sinapsis y en el núcleo del rafe, donde casi no hay sinapsis; no depende del bloqueo de la recaptación ni de los receptores; hay muchas menos moléculas por terminal que transportadores y receptores, y la concentración extracelular coincide con la afinidad del receptor principal. En el núcleo del rafe la inhibición a través de 5-HT$_{1A}$ es punto a punto: el transportador impide que la serotonina se acumule fuera de la hendidura. & \cite{Bunin1998, CourtneyFord2016} & medido \\
6 & La concentración según~\eqref{eq:seroneuron} decae por la recaptación; $\tau_c$ del orden de un segundo es una estimación & Voltamperometría en el cerebro vivo: la serotonina extracelular la limitan ante todo la recaptación y la degradación, no la síntesis. En los trabajos citados no hay medida directa de $\tau_c$; el orden de magnitud es una estimación del capítulo. & \cite{Hashemi2012serotonin} & medido; $\tau_c$, estimación \\
7 & NO y NO-O con un solo marcapasos; completitud de la lógica \nt{SERO} (proposición~\ref{prop:serocomplete}, fig.~\ref{fig:serogates}) & Se sigue de~\eqref{eq:seroneuron}: un marcapasos inhibible es un NO-O, y el NO-O es funcionalmente completo. No hay experimentos en que se haya construido lógica con neuronas \nt{SERO}; la condición de volúmenes separados no se cumple en el cerebro. & deducción en el capítulo & teoría \\
8 & La serotonina inhibe a la propia neurona \nt{SERO} a través de receptores situados en ella & En ratas anestesiadas, los agonistas de 5-HT$_{1A}$ por vía intravenosa y por iontoforesis inhiben la descarga de las neuronas del rafe dorsal con la misma relación dosis-respuesta que la propia serotonina; los agonistas de 5-HT$_{1B}$ casi no actúan. En la rebanada, la serotonina endógena de las células vecinas produce a través de 5-HT$_{1A}$ una corriente y una pausa breve en la descarga, sin cambiar la frecuencia media. & \cite{Sprouse1987, CourtneyFord2016} & medido \\
9 & En el modelo, el lazo a través de un volumen común es un regulador de nivel: la dispersión y la constante de tiempo se reducen $1+G$ veces~\eqref{eq:feedback}; que el lazo funcione así en el cerebro es una hipótesis & Fórmula clásica del amplificador con realimentación negativa; en el capítulo se deduce de nuevo. La ganancia de lazo $G$ del sistema \nt{SERO} no se ha medido. Medido: en la rebanada la autoinhibición es punto a punto, da una pausa breve y no cambia la frecuencia media. & \cite{Black1934feedback, CourtneyFord2016} & teoría; en el cerebro, hipótesis \\
10 & La concentración es una media móvil de la frecuencia, un filtro paso bajo~\eqref{eq:lowpass}, fig.~\ref{fig:lowpass} & Solución de la ecuación lineal~\eqref{eq:seroneuron}; la figura es un cálculo con esta fórmula para $\tau_c=1\,\text{s}$. No hay un modelo aparte en \texttt{models/}. & deducción en el capítulo & teoría \\
11 & La tortuosidad de los espacios frena la difusión de una molécula pequeña unas $2.6$ veces & Método de fuente puntual: iontoforesis de tetrametilamonio y registro de su concentración con un electrodo ionoselectivo en rebanadas y en el cerebro vivo. Tortuosidad $\lambda\approx1.6$, es decir, el $D$ efectivo es $\lambda^2\approx2.6$ veces menor que el libre; fracción de volumen extracelular, cerca de $0.2$. El coeficiente de difusión de la propia serotonina en el tejido no se midió en estos trabajos; en el capítulo es una estimación. & \cite{Sykova2008} & medido \\
12 & Longitud de un “cable” independiente $\ell\sim\sqrt{D\tau}\approx20\um$~\eqref{eq:diffusionlength}; el número de cables está limitado por el número de tales regiones (proposición~\ref{prop:volumewires}) & Ley de difusión y sustitución de las estimaciones de $D$ y $\tau$ (filas~6 y~11); la proposición~\ref{prop:volumewires} es una consecuencia. No hay experimentos directos que cuenten los canales independientes de la transmisión de volumen. & deducción en el capítulo & teoría \\
13 & La salida de una neurona \nt{SERO} se reparte por enormes zonas del cerebro; se estiman $\sim10^5$ sitios de liberación & Reconstrucción completa de neuronas individuales del rafe dorsal en la rata: todos los axones ascendentes se ramifican mucho, la longitud axonal total de una célula llega a $18.7$~cm, y las ramas de una misma célula alcanzan el estriado, la corteza prefrontal y la amígdala. El número de sitios de liberación por célula no se ha contado directamente. & \cite{GagnonParent2014} & medido; $10^5$, estimación \\
14 & Las neuronas \nt{SERO} se dividen en varios subsistemas con salidas y respuestas distintas & Marcaje viral-genético en ratones: las neuronas que van a la amígdala y a la corteza frontal casi no se solapan en su ramificación, reciben entradas distintas y responden de forma opuesta a un estímulo desagradable; activar y desactivar cada grupo cambia la conducta de manera diferente. & \cite{Ren2018} & medido \\
15 & Condición de paciencia $D<\tau_\gamma\ln(R/r_0)$~\eqref{eq:patience} con descuento exponencial; con el medido, hiperbólico, $D<\tau_\gamma(R/r_0-1)$ & Se sigue del descuento $e^{-D/\tau_\gamma}$ o $1/(1+D/\tau_\gamma)$; deducción en el capítulo. Monos, elección con esperas de segundos: el descuento se parece más a una hipérbola que a una exponencial; la respuesta de las neuronas \nt{DOPA} a la señal que anuncia la recompensa diferida cae con la espera también como una hipérbola. & deducción en el capítulo; \cite{KobayashiSchultz2008} & teoría \\
16 & \nt{SERO} fija el horizonte $\tau_\gamma$ (sección~\ref{sec:horizon}) & La activación optogenética de las neuronas \nt{SERO} del rafe dorsal en ratones alarga la espera de una recompensa diferida y aumenta la persistencia en buscar una recompensa que se ha vuelto escasa. En humanos, bajar la serotonina (dieta sin triptófano) hace más empinado el descuento de la recompensa diferida. Cómo se convierte la concentración en $\tau_\gamma$ no se sabe. & \cite{Doya2002, Miyazaki2014, Lottem2018, Schweighofer2008} & hipótesis \\
\bottomrule
\end{longtable}}

\section{Capítulo~\ref{sec:dofa}. Aprendizaje con \nt{DOPA}}

Los mecanismos de la sinapsis (porciones, detector de coincidencias, calcio, ventanas de plasticidad) y el comportamiento de la dopamina como error de predicción están medidos directamente; que las reglas lleven al gradiente y lo que cuesta la difusión son teoremas; el resultado del aprendizaje de elección es un cálculo con el modelo del libro; el circuito que calcula el error es una hipótesis.

{\footnotesize
\setlength{\tabcolsep}{3pt}\renewcommand{\arraystretch}{1.15}
\begin{longtable}{>{\raggedright}p{0.5cm}>{\raggedright}p{4.0cm}>{\raggedright}p{5.6cm}>{\raggedright}p{2.3cm}>{\raggedright\arraybackslash}p{1.7cm}}
\toprule
N.º & Afirmación & Experimento y qué se midió & Fuente & Estado \\
\midrule\endfirsthead
\toprule
N.º & Afirmación & Experimento y qué se midió & Fuente & Estado \\
\midrule\endhead
1 & En el cerebro no se ha hallado retropropagación del error en la forma de las redes artificiales & No hay experimento directo: es una afirmación de ausencia. Las revisiones analizan qué exige el algoritmo (conexiones de vuelta que repitan las de ida, una fase aparte) y qué aproximaciones biológicas se han propuesto. & \cite{Lillicrap2020, WhittingtonBogacz2019} & hipótesis \\
2 & El peso se descompone en número de sitios de liberación, probabilidad de liberación y respuesta a una porción: $w=N_s\,p\,A_1$~\eqref{eq:quantal} & Unión neuromuscular de rana con liberación reducida: la respuesta del destinatario varía en escalones múltiplos de la respuesta miniatura espontánea a una porción; el número de porciones por espiga sigue la ley de Poisson. & \cite{delCastillo1954} & medido \\
3 & El receptor del destinatario es un detector de coincidencias; el calcio dispara el cambio de peso & Patch-clamp en neuronas de ratón: los iones Mg$^{2+}$ taponan el canal abierto por el glutamato en reposo y salen con la despolarización. Rebanadas de hipocampo: un quelante de calcio en el destinatario bloquea la potenciación a largo plazo, y liberar calcio con luz dentro de la célula refuerza por sí solo la transmisión. & \cite{Nowak1984, Malenka1988calcium} & medido \\
4 & La decisión la ejecuta cualquiera de los lados: crece $A_1$ en el destinatario, cambia $p$ en el emisor & El glutamato liberado con luz sobre una sola espina provoca un crecimiento rápido y duradero de la espina y de la corriente por sus receptores; la potenciación va acompañada de la inserción de receptores AMPA. La despolarización del destinatario suprime temporalmente la liberación en el emisor; el efecto lo transmiten endocannabinoides que cruzan la hendidura hacia atrás (mostrado en entradas \nt{GABA}). La depresión a corto plazo depende de la probabilidad de liberación en el emisor. & \cite{Matsuzaki2004, Malinow2002, WilsonNicoll2001, Tsodyks1997} & medido \\
5 & La sinapsis se refuerza cuando el emisor y el destinatario están activos a la vez~\eqref{eq:hebb} & Conejo anestesiado: tras series breves de espigas por la vía de entrada del hipocampo, la respuesta queda reforzada de $30$~min a $10$~h. En una rebanada de hipocampo, las espigas del destinatario refuerzan solo las sinapsis activas en ese mismo momento. & \cite{Bliss1973LTP, Kelso1986hebb} & medido \\
6 & La regla~\eqref{eq:oja} halla el vector propio principal de las correlaciones de las entradas (proposición~\ref{prop:oja}) & Teorema; demostración en el capítulo. No hay experimento en que una neurona haya aprendido la componente principal de sus entradas; el mecanismo que limita el crecimiento de los pesos en la sinapsis no se conoce. & \cite{Oja1982} & teoría \\
7 & Hebb temporal: ventana de unos $\pm20\ms$ (fig.~\ref{fig:stdp}); de secuencias repetidas crecen cadenas & Pares de neuronas de hipocampo en cultivo: una espiga del destinatario hasta $20\ms$ después de la del emisor da un refuerzo duradero; hasta $20\ms$ antes, un debilitamiento; ambos dependen de los receptores NMDA. La relación $A_-=0.6A_+$ de la figura es ilustrativa. Que así crezcan cadenas en la red no se ha medido directamente. & \cite{BiPoo1998} & medido \\
8 & Traza~\eqref{eq:threefactor}: la dopamina que llega $1$--$2\,\text{s}$ después de la actividad conjunta refuerza la conexión; más tarde, no (proposición~\ref{prop:threefactor}) & Rebanadas de estriado, estimulación optogenética separada de las entradas glutamatérgicas y dopaminérgicas: la espina crece solo si la dopamina llega $0.3$--$2\,\text{s}$ después de la entrada glutamatérgica; la ventana la fija la subida y caída rápidas del AMPc. En la corteza, el apareamiento deja trazas separadas para el refuerzo y el debilitamiento, que los receptores de monoaminas convierten en cambios duraderos en una ventana del orden de segundos. & \cite{Yagishita2014, He2015eligibility} & medido \\
9 & También hay ventanas de segundos sin dopamina: la tercera señal es una excitación fuerte del destinatario & Ratón vivo, campo CA1: un solo evento de meseta en el destinatario refuerza las entradas llegadas segundos antes y después, y crea un campo de lugar en una sola pasada. & \cite{Bittner2017} & medido \\
10 & El paso medio de la regla de tres factores sigue el gradiente de la recompensa~\eqref{eq:perturbation} (proposición~\ref{prop:gradient}) & Teorema del gradiente para el aprendizaje por perturbaciones aleatorias; en el capítulo se deduce para ruido pequeño. & \cite{Williams1992} & teoría \\
11 & El ruido es búsqueda: la red aprende de sus propias desviaciones aleatorias & Diamantes mandarín jóvenes: desactivar el núcleo de salida del circuito de los ganglios basales elimina de forma brusca y reversible la variabilidad del canto. Diamantes adultos: se aplica una perturbación a las ejecuciones de una sílaba cuyo tono queda a un lado del medio, y las aves desplazan pronto el tono hacia el otro lado: aprenden de su propia dispersión. & \cite{Olveczky2005, TumerBrainard2007} & medido \\
12 & La elección entre dos se aprende con $\tau_e=1\,\text{s}$ y $\delta=R-\bar R$; no se aprende con $\tau_e=50\ms$; sin restar, los pesos crecen sin límite, y con una tasa de aprendizaje alta la red puede fijar la peor elección & \texttt{models/dofa\_learning.py}, $\eta=0.05$, $500$ intentos, $6$ semillas. $\tau_e=1\,\text{s}$: fracción de elecciones de $A$ $0.65$--$0.79\to0.96$--$1.00$, pesos $0.85$--$1.33$. $\tau_e=50\ms$: $0.38$--$0.55$, los pesos no cambian. $\delta=R$: peso del ganador $860$--$1190$. $\delta=R$, $\eta=0.5$, $3$ semillas: una se fijó en $B$ ($0.04\to0$). El script imprime los cuatro casos; el recálculo coincidió con el capítulo. & \texttt{models/\allowbreak dofa\_\allowbreak learning.py} & modelo \\
13 & El tiempo de aprendizaje con una única señal común crece en proporción al número de neuronas ruidosas (proposición~\ref{prop:broadcastcost}) & Curvas de aprendizaje de una red lineal: la perturbación de neuronas es más lenta que el gradiente exacto tantas veces como neuronas de salida hay, y la perturbación de pesos, además, tantas veces como entradas. & \cite{Werfel2005} & teoría \\
14 & Las salidas de \nt{DOPA} van sobre todo a las regiones de elección de acciones & Reconstrucción completa de los axones de neuronas \nt{DOPA} nigroestriatales individuales de rata: las ramas de una célula cubren del $0.45$ al $5.7\%$ del volumen del estriado (en promedio, el $2.7\%$); fuera del estriado casi no hay ramas. & \cite{Matsuda2009} & medido \\
15 & La corteza aprende a predecir su entrada, y el error se calcula en el sitio & Revisión: en la corteza sensorial hay respuestas al desajuste entre la entrada esperada y la recibida. Que sea una regla general de aprendizaje de la corteza no está demostrado. & \cite{Keller2018} & hipótesis \\
16 & La dopamina se comporta como el error~\eqref{eq:ctd}: pico, traslado a la señal previa, valle (fig.~\ref{fig:ctdcases}) & Neuronas \nt{DOPA} de mono: pico ante un jugo inesperado; tras el aprendizaje, pico ante la señal previa y ninguna respuesta al jugo prometido; valle cuando el jugo no llega. La forma continua~\eqref{eq:ctd} es teoría. & \cite{Schultz1997, Doya2000} & medido \\
17 & Circuito que calcula $\dd V/\dd t$ y resta la expectativa & Ratones, registro y optogenética en el área tegmental ventral: las neuronas \nt{DOPA} restan la expectativa de la recompensa; las neuronas \nt{GABA} vecinas las inhiben cuando se espera la recompensa, y excitarlas o inhibirlas cambia el error. Cómo se calcula la derivada no se ha mostrado. & \cite{Eshel2015arithmetic} & hipótesis \\
18 & La neurona \nt{DOPA} es un marcapasos; los valles son menos profundos que los picos & En la rebanada, las neuronas \nt{DOPA} descargan de forma espontánea y muy regular, sobre un potencial marcapasos lento. En el mono, la frecuencia sigue cuantitativamente el error positivo, y el negativo lo transmite solo débilmente. & \cite{GraceOnn1989, Bayer2005} & medido \\
19 & Actor y crítico (fig.~\ref{fig:actorcritic}) & El algoritmo procede de la teoría del aprendizaje por refuerzo. En humanos (RMf), el estriado ventral se comporta como crítico con elección y sin ella; el dorsal, solo cuando hay que elegir acciones. & \cite{SuttonBarto2018, ODoherty2004actor} & hipótesis \\
20 & El signo de $\delta$ en la regla está invertido en las neuronas con la segunda familia de receptores & Rebanadas de estriado: en las neuronas con receptores D1 el apareamiento da un refuerzo que requiere D1; en las neuronas con D2, un debilitamiento que requiere D2. & \cite{Shen2008} & medido \\
\bottomrule
\end{longtable}}

\section{Capítulo~\ref{sec:dofaalt}. Otras teorías de \nt{DOPA}}

Cada uno de los cuadros ampliados se apoya en sus propias medidas directas de la dopamina (dispersión de los puntos de inversión, respuestas a lo desagradable y a lo nuevo, crecimiento lento, respuestas al cambio de sabor), pero las fórmulas que las explican son teorías, y entre dos de ellas los experimentos se contradicen por ahora.

{\footnotesize
\setlength{\tabcolsep}{3pt}\renewcommand{\arraystretch}{1.15}
\begin{longtable}{>{\raggedright}p{0.5cm}>{\raggedright}p{4.0cm}>{\raggedright}p{5.6cm}>{\raggedright}p{2.3cm}>{\raggedright\arraybackslash}p{1.7cm}}
\toprule
N.º & Proposición & Experimento y qué se midió & Fuente & Estado \\
\midrule\endfirsthead
\toprule
N.º & Proposición & Experimento y qué se midió & Fuente & Estado \\
\midrule\endhead
1 & La regla asimétrica~\eqref{eq:asymmetric} converge al punto~\eqref{eq:expectile}; con $\kappa=1/2$ es la media; para una gota con probabilidad $1/2$, $V=\kappa$ (fig.~\ref{fig:expectiles}) & El punto~\eqref{eq:expectile} es un expectil, solución de un problema de mínimos cuadrados con pesos distintos para las desviaciones positivas y negativas; para el ejemplo del capítulo, la deducción está en el propio capítulo. & \cite{NeweyPowell1987}; deducción en el capítulo & teoría \\
2 & Distintas neuronas \nt{DOPA} tienen distintos puntos de inversión, ordenados según la asimetría (proposición~\ref{prop:expectile}) & Ratones, neuronas \nt{DOPA} del área tegmental ventral marcadas optogenéticamente, recompensas de distinto tamaño: el punto en que el pico pasa a valle difiere entre neuronas; es más alto en las neuronas con respuestas más asimétricas; a partir de las respuestas del grupo se reconstruye la forma de la distribución de recompensas. Que esta sea la función principal de la dispersión es una hipótesis. & \cite{Dabney2020} & medido \\
3 & Parte de las neuronas \nt{DOPA} responde a lo desagradable con un pico & Monos: unas neuronas \nt{DOPA} se excitan con la señal previa a la recompensa y se inhiben con la señal previa a un soplo de aire en el ojo; otras se excitan con ambas; los grupos están en partes distintas del mesencéfalo. & \cite{Matsumoto2009, BrombergMartin2010} & medido \\
4 & El módulo del error o la novedad fija la tasa de aprendizaje & En los trabajos citados no hay experimento directo en que la señal de importancia de las neuronas \nt{DOPA} cambie la tasa de aprendizaje; la revisión propone el papel de una señal de “atención, es importante”. & \cite{BrombergMartin2010} & hipótesis \\
5 & Un cable aparte para el eje del peligro & Ratones: las neuronas \nt{DOPA} que van al extremo posterior del estriado responden a estímulos nuevos e intensos, no a la recompensa; su destrucción debilita la evitación de un objeto nuevo. & \cite{Menegas2018} & medido \\
6 & Tiempo óptimo de la acción $\tau_a^*=\sqrt{C/\bar R}$~\eqref{eq:vigor} & Mínimo de la suma $C/\tau_a+\bar R\tau_a$; deducción en el capítulo. En el modelo original, la rapidez de las acciones la fija el precio del tiempo, igual a la recompensa media. & \cite{Niv2007}; deducción en el capítulo & teoría \\
7 & El nivel lento de dopamina lleva $\bar R$ y fija la rapidez de las acciones (proposición~\ref{prop:multiplex}) & Ratas, microdiálisis y voltamperometría en el núcleo accumbens: el nivel de dopamina sigue minuto a minuto la frecuencia de recompensas y la rapidez de las acciones. En humanos, la L-DOPA refuerza la dependencia del tiempo de reacción respecto de la frecuencia media de recompensas. Que el nivel lento sea precisamente $\bar R$ no está demostrado. & \cite{Hamid2016, Beierholm2013vigor} & hipótesis \\
8 & El nivel lento se forma a partir de dos fuentes: la liberación cambia en las salidas sin que cambie la frecuencia de espigas, pero también hay crecimiento lento en la propia frecuencia & Ratas, una misma tarea: la liberación en el núcleo accumbens sigue la expectativa cambiante de recompensa, y la frecuencia de espigas de las neuronas \nt{DOPA} identificadas del área tegmental ventral, no. Ratones en un pasillo virtual (fila~10): el crecimiento gradual al acercarse a la recompensa está también en las espigas de neuronas \nt{DOPA} identificadas. & \cite{Mohebi2019, Kim2020} & medido \\
9 & La dopamina crece de forma gradual mientras el animal se acerca a una recompensa lejana & Ratas en un laberinto, voltamperometría en el estriado: la concentración de dopamina sube en segundos a medida que se acercan a la meta; la pendiente depende de la distancia y del tamaño de la recompensa. & \cite{Howe2013ramp, Hamid2016} & medido \\
10 & El crecimiento es el error $\delta$ al afinarse la estimación~\eqref{eq:ramp}; un salto en el pasillo da un pico (fig.~\ref{fig:ramp}) & La fórmula y el número $\delta=0.75\,V$ por segundo son deducción del capítulo. Experimento: ratones en un pasillo virtual, traslado instantáneo más cerca de la meta; las espigas de los somas, el calcio de somas y salidas y la concentración en el estriado responden como un error, no como la expectativa $V$. Cuál de las dos explicaciones vale en todas las salidas se discute. & \cite{Kim2020} & medido \\
11 & Mirar atrás: la dopamina informa de una medida del vínculo causal~\eqref{eq:retro} & Ratones, liberación de dopamina en el núcleo accumbens en experimentos en que esta teoría y el error~\eqref{eq:ctd} predicen cosas distintas: en varios experimentos la liberación siguió a la primera. & \cite{Jeong2022} & hipótesis \\
12 & La respuesta de la dopamina durante el aprendizaje se desplaza poco a poco de la recompensa a la señal previa & Ratones, señal de las salidas de las neuronas \nt{DOPA} en el estriado ventral a lo largo del aprendizaje: el pico pasa gradualmente del momento de la recompensa al de la señal previa, como exige el aprendizaje según~\eqref{eq:ctd}. & \cite{Amo2022} & medido \\
13 & Las neuronas \nt{DOPA} responden al cambio de sabor de la recompensa con el mismo valor & Ratas, registro de neuronas del área tegmental ventral: sustituir un jugo por otro de igual valor provoca respuesta, aunque el error~\eqref{eq:ctd} es cero. & \cite{Takahashi2017} & medido \\
14 & Picos artificiales enseñan el vínculo entre dos estímulos neutros & Ratas, experimento de precondicionamiento sensorial con dos estímulos sin recompensa: la activación optogenética de las neuronas \nt{DOPA} en la transición del primero al segundo crea el vínculo, y su inhibición impide aprenderlo. & \cite{Sharpe2017} & medido \\
15 & Vector de errores por rasgos~\eqref{eq:vectortd} & Teoría del error de predicción de rasgos; no hay comprobación directa de la forma vectorial. & \cite{Gardner2018} & hipótesis \\
16 & Distintas neuronas \nt{DOPA} llevan en una misma tarea magnitudes distintas & Ratones en un pasillo virtual, imagen de dos fotones de neuronas \nt{DOPA}: unas se relacionan con la recompensa, otras con la velocidad y el giro, otras con las señales previas y la precisión de la elección. & \cite{Engelhard2019} & medido \\
17 & Un mazo de cables en lugar de un cable (proposición~\ref{prop:bundle}) & Resumen de las filas~2--16: cada descomposición está medida por separado; no hay un cuadro experimental único que muestre que todas son partes de una misma señal. & --- & hipótesis \\
\bottomrule
\end{longtable}}

\section{Capítulo~\ref{sec:nora}. \nt{NORA}}

La organización del sistema \nt{NORA}, sus dos modos, el vínculo con la pupila (no solo a través de \nt{NORA}), el aumento de la relación señal/fondo y la U invertida en la conducta están medidos directamente; la ganancia multiplicativa $g$ es un modelo; que la ganancia conmute el modo de elección y actúe como temperatura inversa son deducciones del capítulo; el beneficio de ajustar la ganancia es un cálculo con el modelo del libro; “\nt{NORA} es la perilla de la exploración” y “\nt{NORA} es una interrupción” son hipótesis.

{\footnotesize
\setlength{\tabcolsep}{3pt}\renewcommand{\arraystretch}{1.15}
\begin{longtable}{>{\raggedright}p{0.5cm}>{\raggedright}p{4.0cm}>{\raggedright}p{5.6cm}>{\raggedright}p{2.3cm}>{\raggedright\arraybackslash}p{1.7cm}}
\toprule
N.º & Proposición & Experimento y qué se midió & Fuente & Estado \\
\midrule\endfirsthead
\toprule
N.º & Proposición & Experimento y qué se midió & Fuente & Estado \\
\midrule\endhead
1 & En el ser humano hay unas decenas de miles de neuronas \nt{NORA}, casi todas en un solo grupo & Cortes seriados del tronco encefálico humano, tinción de tirosina hidroxilasa: $53\,900$ células en el locus coeruleus y otras $6\,260$ en el subnúcleo vecino. & \cite{Baker1989LC} & medido \\
2 & Las salidas se extienden por casi todo el cerebro, pero partes del grupo atienden en parte a regiones distintas & Marcaje viral-genético en ratones: las neuronas \nt{NORA} del locus coeruleus reciben entradas de amplias regiones del cerebro y envían salidas a amplias regiones del cerebro y de la médula espinal. En ratas: el grupo que va a la amígdala refuerza el aprendizaje del miedo, y un grupo aparte que va a la corteza prefrontal, su extinción. & \cite{Schwarz2015LC, Uematsu2017LC, Poe2020} & medido \\
3 & Modo de fondo: espigas regulares con una frecuencia de fracciones de espiga a unas pocas por segundo, que cambia con la vigilia & Ratas, registro de neuronas del locus coeruleus a lo largo del ciclo sueño-vigilia: la frecuencia es máxima en vigilia, menor en el sueño lento y casi nula en el sueño REM; los cambios de frecuencia se adelantan al cambio de estado. & \cite{AstonJonesBloom1981} & medido \\
4 & Modo fásico: una ráfaga breve ante un suceso importante para la tarea, aproximadamente en el momento de decidir & Monos, tarea de detectar un blanco poco frecuente: todas las neuronas del locus coeruleus responden con una ráfaga solo al blanco, $\sim90\ms$ después de él y $\sim200\ms$ antes de la respuesta manual; cuando el desempeño es malo, las ráfagas son más débiles. En una tarea de elección entre dos, la ráfaga está más ligada a la respuesta que al estímulo, y aparece tanto en respuestas correctas como erróneas. & \cite{AstonJones1994vigilance, Clayton2004LC} & medido \\
5 & El diámetro de la pupila es un punto de control impreciso de \nt{NORA}: sigue a \nt{NORA}, a \nt{ACET} y a otras regiones & Monos: las espigas del locus coeruleus, hasta espigas sueltas de una célula, predicen la dilatación de la pupila; hay un vínculo parecido pero más tardío en los colículos y la corteza cingulada. Ratones: las fluctuaciones de la pupila siguen la actividad de las salidas tanto noradrenérgicas como colinérgicas en la corteza. & \cite{Joshi2016, Reimer2016} & medido \\
6 & La noradrenalina suprime la actividad de fondo más que la evocada; la ganancia multiplicativa~\eqref{eq:gain} es un modelo que reproduce el efecto & Monos despiertos, corteza auditiva, iontoforesis de noradrenalina: la actividad de fondo de las neuronas se suprime más que la respuesta a las vocalizaciones de congéneres; la relación señal-ruido crece. La sigmoide multiplicativa~\eqref{eq:gain} procede de un modelo de red; la multiplicación literal de la pendiente no se ha medido. & \cite{Foote1975NE, ServanSchreiber1990} & medido; $g$, modelo \\
7 & La ganancia aumenta $d'$ solo frente al ruido posterior al amplificador~\eqref{eq:gainsnr} (proposición~\ref{prop:gainnoise}) & Deducción en el capítulo; en el modelo de red, la ganancia mejora la detección de la señal. No hay experimento que separe el ruido anterior y posterior al amplificador y compruebe esta diferencia. & deducción en el capítulo; \cite{ServanSchreiber1990} & teoría \\
8 & La frontera $g\beta=1$ conmuta la red de elección de “delibera” a “ha decidido”~\eqref{eq:wtagain} (proposición~\ref{prop:gainwta}, fig.~\ref{fig:pitchfork}) & Análisis lineal de dos grupos que se inhiben mutuamente; deducción en el capítulo. No hay medidas directas de este umbral en el cerebro. & deducción en el capítulo & teoría \\
9 & La ganancia es la temperatura inversa de la elección~\eqref{eq:softmax} & Con ruido logístico después del amplificador, la probabilidad de elección es un softmax con $T=\sigma/g$; deducción en el capítulo. & deducción en el capítulo & teoría \\
10 & \nt{NORA} fija cuánto explorar: una actividad de fondo alta es una $g$ efectiva pequeña (exploración); la ráfaga fásica en el momento de decidir, una grande (decisión) & Humanos: antes de una elección al azar la pupila basal es más ancha que antes de elegir la mejor opción conocida; las personas con pupila más ancha tienden más a explorar. Ratas: el paso al modo de elección aleatoria lo controla la entrada de \nt{NORA} a la corteza cingulada anterior. Que la ráfaga aumente la $g$ efectiva no se ha medido directamente. & \cite{JepmaNieuwenhuis2011, Tervo2014stochastic, AstonJones2005} & hipótesis \\
11 & La recompensa depende de $g$ como una U invertida; la mejor $g$ es menor en un mundo que cambia deprisa (proposición~\ref{prop:explore}, fig.~\ref{fig:invertedU}) & \texttt{models/nora\_explore.py}, $60$ ejecuciones de $4000$ intentos. Un cambio cada $1000$ intentos: mejor $g=16$, fracción $0.976$; cada $20$: $g=8$, fracción $0.886$; con $g=0.5$, $0.77$. El recálculo coincidió con el capítulo. & \texttt{models/\allowbreak nora\_\allowbreak explore.py} & modelo \\
12 & U invertida en la conducta: el mejor desempeño, con actividad de fondo moderada y ráfagas nítidas & Monos, la misma tarea que en la fila~4: en los periodos de buen desempeño la frecuencia de fondo es moderada y las ráfagas ante el blanco son nítidas; con frecuencia de fondo alta no hay ráfagas y hay más falsas alarmas. Los cambios de la actividad de fondo y evocada siguen las fluctuaciones del desempeño. & \cite{Usher1999LC, AstonJones2005} & medido \\
13 & \nt{NORA} comunica la incertidumbre inesperada, y \nt{ACET}, la esperada & Teoría de inferencia bayesiana con dos clases de incertidumbre; en los trabajos citados no hay experimento directo que separe estas señales por neurotransmisor. & \cite{YuDayan2005} & hipótesis \\
14 & Tras un cambio brusco las personas aprenden más deprisa, y la pupila se dilata & Humanos, tarea de predicción con cambios repentinos: los cambios breves de la pupila siguen la estimación de la probabilidad de cambio y, junto con la pupila basal, predicen cuánto modifican la estimación los datos nuevos (tasa de aprendizaje); un cambio de pupila ajeno a la luz y a la tarea modifica esa tasa. Que aquí la pupila indique precisamente \nt{NORA} es una inferencia a partir de los datos indirectos de la fila~5. & \cite{Nassar2012} & medido \\
15 & La ráfaga fásica funciona como una interrupción: la red borra su estado & No hay experimento directo en que una ráfaga de \nt{NORA} borre el estado de la red; solo están medidas las ráfagas ante sucesos importantes (fila~4). & \cite{BouretSara2005} & hipótesis \\
16 & Ajustar $g$ o la tasa de aprendizaje según la sorpresa apenas mejora los mejores ajustes constantes & \texttt{models/nora\_explore.py}, mundo con épocas de $1000$ intentos (un cambio cada $1000$ y cada $20$). Mejor constante, $g=8$: $0.925$; ajuste de $g$: hasta $0.927$; ajuste de la tasa de aprendizaje: hasta $0.926$; tasa constante $0.4$: $0.928$. El recálculo coincidió con el capítulo. & \texttt{models/\allowbreak nora\_\allowbreak explore.py} & modelo \\
\bottomrule
\end{longtable}}

\section{Capítulo~\ref{sec:acet}. Añadimos \nt{ACET}}

Las tres acciones de la acetilcolina están medidas en rebanadas y en el cerebro vivo; el conflicto entre escritura y lectura se muestra con el modelo del libro; y que \nt{ACET} sirva de línea de habilitación de escritura y comunique la incertidumbre esperada es una hipótesis con apoyo experimental parcial.

{\footnotesize
\setlength{\tabcolsep}{3pt}\renewcommand{\arraystretch}{1.15}
\begin{longtable}{>{\raggedright}p{0.5cm}>{\raggedright}p{4.0cm}>{\raggedright}p{5.6cm}>{\raggedright}p{2.3cm}>{\raggedright\arraybackslash}p{1.7cm}}
\toprule
N.º & Afirmación & Experimento y qué se midió & Fuente & Estado \\
\midrule\endfirsthead
\toprule
N.º & Afirmación & Experimento y qué se midió & Fuente & Estado \\
\midrule\endhead
1 & Las neuronas \nt{ACET} del cerebro son pocas, están reunidas en unos cuantos grupos, y sus salidas se extienden por toda la corteza: un canal de difusión & Tinción con anticuerpos contra la enzima de síntesis de la acetilcolina y marcaje retrógrado en la rata: la corteza, el hipocampo, la amígdala, el bulbo olfatorio y el tálamo reciben acetilcolina sobre todo de seis grupos compactos de células del prosencéfalo basal y del tronco & \cite{Mesulam1983} & medido \\
2 & El nivel de acetilcolina en la corteza es alto en la vigilia y bajo en el sueño lento & Microdiálisis en la corteza y el hipocampo de gatos en libertad de movimiento con registro de EEG: la liberación de acetilcolina en vigilia tranquila es $\sim100\%$ mayor, y en vigilia activa $\sim175\%$ mayor, que en el sueño lento; en el sueño REM, en la corteza, como en vigilia & \cite{Marrosu1995,Hasselmo1999} & medido \\
3 & El sentido de la señal lo fija el receptor: la acetilcolina tiene receptores-canal rápidos y receptores lentos que disparan reacciones en la célula (proposición~\ref{prop:receptor}) & Revisión de medidas: el efecto de la acetilcolina sobre la excitabilidad, la transmisión y la plasticidad depende del lugar de liberación, del subtipo de receptor (los nicotínicos son canales iónicos; los muscarínicos actúan a través de proteínas G) y de la célula destinataria & \cite{Picciotto2012} & medido \\
4 & Primera acción: debilitar las conexiones internas casi sin tocar las entradas externas (factor $1-a$ en~\eqref{eq:acetnet}) & Rebanadas de corteza olfativa de rata: con $100$\,\textmu M de agonistas de la acetilcolina, los potenciales de las conexiones internas (capa Ib) caen más de un $60\%$, y los de la entrada externa (capa Ia), menos de un $15\%$, en la misma célula; la supresión es presináptica. Rebanadas de hipocampo (CA1): $89\%$ frente a $40\%$ para las vías interna y de entrada & \cite{HasselmoBower1992,HasselmoSchnell1994} & medido \\
5 & Segunda acción: reforzar la entrada (factor $\frac{1+a}{2}$) & Rebanadas de neocorteza: los receptores nicotínicos refuerzan las sinapsis de la entrada talámica y no actúan sobre las intracorticales; los muscarínicos debilitan ambos tipos. La acetilcolina aumenta la excitabilidad de las neuronas (revisión) & \cite{Gil1997,Hasselmo2006} & medido \\
6 & Tercera acción: facilitar el cambio de los pesos (tasa $\eta a$) & Rata: la estimulación eléctrica del núcleo basal (la fuente de acetilcolina de la corteza) emparejada con un tono provoca en el animal adulto una fuerte reorganización del mapa de la corteza auditiva hacia el sonido presentado & \cite{KilgardMerzenich1998,Hasselmo2006} & medido \\
7 & La regla de Hebb para patrones dispersos de la segunda ecuación~\eqref{eq:acetnet} da una memoria asociativa que completa el patrón a partir de una parte & Cálculo de la capacidad de una red con patrones dispersos y regla de covarianza: con una fracción pequeña $p$ de neuronas activas, la red guarda bastantes más patrones que con $p=1/2$ & \cite{Tsodyks1988} & teoría \\
8 & Escribir con $a$ bajo daña la memoria: la red escribe lo recordado, no lo visto; con $a=0.1$ el daño no es menor que con $a=0$ & \texttt{models/\allowbreak acet\_memory.py}: $200$ neuronas, cinco patrones de $20$, el patrón nuevo comparte $10$ neuronas con uno de ellos, $10$ ejecuciones. Con $a=0$ el patrón nuevo no se memoriza, y el viejo arrastra $5.9$ neuronas ajenas de $10$; con $a=0.1$ el nuevo se recuerda ($0.97$), pero ambos se funden: el nuevo arrastra $7.3$ neuronas ajenas y el viejo, $7.9$; desde $a=0.2$, menos de una & deducción en el capítulo & modelo \\
9 & Proposición~\ref{prop:writeread}: no hay un nivel $a$ con el que escritura y lectura sean igual de buenas (fig.~\ref{fig:acetsim}) & El mismo script, tasa de aprendizaje $\eta a$: el patrón nuevo se memoriza entero en una presentación con $a\ge0.9$, en $0.8$ con $a=0.75$, y no se memoriza con $a\le0.5$. Lectura a partir de la mitad del patrón: $1.0$ con $a\le0.2$, $0.96$ con $a=0.3$, $0.04$ con $a=0.4$ & deducción en el capítulo & modelo \\
10 & En el cerebro, un solo cable de difusión, \nt{ACET}, conmuta escritura y lectura & La idea procede de modelos construidos a partir de medidas en rebanadas. El experimento más directo es en humanos: el bloqueante de receptores muscarínicos escopolamina empeora el aprendizaje de pares de palabras nuevos, sobre todo de los que más se solapan con los viejos, pero no el recuerdo de pares ya aprendidos. No hay medida directa de los dos modos de la red & \cite{HasselmoAndersonBower1992,Atri2004} & hipótesis \\
11 & El filtro~\eqref{eq:kalman} es óptimo; el peso de la entrada y la tasa de aprendizaje son una misma magnitud $P/R$; en régimen estacionario $P=\sqrt{qR}$ y la memoria es $\sqrt{R/q}$ (proposición~\ref{prop:kalman}) & No requiere experimento: la optimalidad del filtro de Kalman--Bucy es un resultado clásico de la teoría de estimación; el régimen estacionario se deduce en el capítulo & \cite{KalmanBucy1961} & teoría \\
12 & \nt{ACET} comunica a la red una magnitud parecida a $P$: la incertidumbre esperada & Propuesto en un modelo probabilístico de la atención. No hay medida directa de que el nivel de acetilcolina siga la incertidumbre de la estimación & \cite{YuDayan2005} & hipótesis \\
13 & Parte de las neuronas \nt{ACET} responden a la recompensa y al castigo en un par de decenas de milisegundos, con más fuerza cuanto más inesperado es el refuerzo & Ratón, neuronas colinérgicas del prosencéfalo basal identificadas optogenéticamente: respuesta a la recompensa y al castigo a los $18\pm3$\,ms; la magnitud crece con lo inesperado del refuerzo; las respuestas son parecidas en neuronas de dos núcleos & \cite{Hangya2015} & medido \\
14 & \nt{NORA} nota un cambio inesperado del mundo, y \nt{ACET} mantiene la red en modo de escritura & El reparto del trabajo se propone en un modelo. Indirectamente: en humanos, el diámetro de la pupila, ligado a \nt{NORA}, sigue los cambios y la tasa de aprendizaje. No hay comprobación directa del par \nt{NORA}--\nt{ACET} & \cite{YuDayan2005,Nassar2012} & hipótesis \\
15 & En el sueño lento la red, en modo de lectura, reproduce lo escrito durante el día, y estas repeticiones lo copian en la memoria a largo plazo & Registros de conjuntos de neuronas del hipocampo de rata: las células que trabajaron juntas de día se disparan juntas con más frecuencia en el sueño lento posterior, y en el mismo orden: medido. Para la reescritura: suprimir las ondas agudas del hipocampo tras el aprendizaje empeora la memoria espacial, y alargar las ondas espontáneas la mejora; que la causa sea el nivel bajo de \nt{ACET} no está demostrado & \cite{WilsonMcNaughton1994,Skaggs1996,Girardeau2009,FernandezRuiz2019,Hasselmo1999} & hipótesis \\
\bottomrule
\end{longtable}}

\section{Capítulo~\ref{sec:synthesis}. La máquina completa}

El capítulo de síntesis no presenta experimentos nuevos: cada fila se apoya en el capítulo donde se analiza la proposición y en las mismas fuentes; el bus \nt{GLUT} y el control del flujo mediante \nt{GABA} están firmemente establecidos, mientras que los papeles de los canales de difusión y su trabajo conjunto son sobre todo una hipótesis.

{\footnotesize
\setlength{\tabcolsep}{3pt}\renewcommand{\arraystretch}{1.15}
\begin{longtable}{>{\raggedright}p{0.5cm}>{\raggedright}p{4.0cm}>{\raggedright}p{5.6cm}>{\raggedright}p{2.3cm}>{\raggedright\arraybackslash}p{1.7cm}}
\toprule
N.º & Afirmación & Experimento y qué se midió & Fuente & Estado \\
\midrule\endfirsthead
\toprule
N.º & Afirmación & Experimento y qué se midió & Fuente & Estado \\
\midrule\endhead
1 & Para $8.6\cdot10^{10}$ neuronas hay solo cuatro clases de señales de control~\eqref{eq:machine} & Recuento de núcleos celulares en un homogeneizado de cerebro (fraccionador isotrópico): un hombre adulto tiene $86.1\pm8.1$ mil millones de neuronas & \cite{Azevedo2009} & medido \\
2 & Proposición~\ref{prop:architecture}, las dos primeras capas: los datos viajan por el bus de direcciones \nt{GLUT}, el signo de la conexión lo fija el emisor, y las neuronas \nt{GABA} guían y normalizan el flujo (capítulos~\ref{sec:neurontypes}, \ref{sec:gaba}) & Un neurotransmisor principal por neurona (revisión de mediciones); la inhibición prealimentada estrecha la ventana en la que la neurona piramidal suma sus entradas; la división por la actividad total se ha medido en muchas áreas & \cite{Strata1999,Pouille2001,Carandini2012} & medido \\
3 & Los elementos no son fiables: la sinapsis pierde la mayoría de las espigas (tabla~\ref{tab:computer}, capítulo~\ref{sec:code}) & Cortes de hipocampo, estimulación mínima: más de la mitad de las espigas no provoca respuesta en CA1; se descartan fallos del umbral y de la conducción por el axón; la causa es la liberación probabilística del neurotransmisor & \cite{AllenStevens1994} & medido \\
4 & El cerebro funciona con $\sim20$ vatios, en promedio cerca de una espiga por segundo por neurona (capítulo~\ref{sec:code}); los ingenieros aún no han construido una máquina así & Estimación~\eqref{eq:energy} a partir de los costos medidos: unas $10^9$ moléculas de ATP por espiga, incluido el trabajo de las sinapsis (corteza de roedores); una estimación detallada para la corteza da una frecuencia aún menor. Estado de la técnica, según una revisión & \cite{Attwell2001,Lennie2003,MehonicKenyon2022} & teoría \\
5 & El aprendizaje de la mayoría de las sinapsis es el producto de una traza local por un único número común $\delta$ (segunda ecuación de~\eqref{eq:machine}, capítulo~\ref{sec:dofa}) & Las neuronas \nt{DOPA} del mono responden al error de predicción de la recompensa; en cortes de estriado, la dopamina agranda la espina solo si llega $0.3$--$2$\,s después de la entrada glutamatérgica: la traza está medida & \cite{Schultz1997,Yagishita2014} & medido \\
6 & Un cable de error propio para cada salida existe solo en el circuito del aprendizaje motor (capítulo~\ref{sec:motorlearning}) & Cerebelo: la coincidencia de la señal de la fibra trepadora con una entrada debilita esa entrada; en el mono, las espigas complejas de las células de Purkinje llevan la dirección del error del movimiento y provocan la corrección & \cite{Ito2001,Herzfeld2018} & medido \\
7 & Cada canal de difusión es un mazo de unas pocas líneas (proposición~\ref{prop:bundle}) & Para \nt{DOPA}: en una misma tarea, distintas neuronas llevan la recompensa, el movimiento y los rasgos del estímulo; el error rápido y el nivel lento de dopamina son distintos. Para \nt{SERO}, \nt{NORA} y \nt{ACET} este desglose está menos estudiado & \cite{Engelhard2019,Hamid2016,Mohebi2019} & medido \\
8 & \nt{SERO}: regulador de nivel y, según la hipótesis, horizonte $\tau_\gamma$ (capítulo~\ref{sec:serocomputer}) & Autoinhibición: los agonistas de sus propios receptores 5-HT$_{1A}$ inhiben las neuronas serotoninérgicas del rafe (medido). El horizonte se propuso en la teoría; el experimento más fuerte: la activación optogenética de estas neuronas en el ratón alarga la espera de una recompensa diferida & \cite{Sprouse1987,Doya2002,Miyazaki2014} & hipótesis \\
9 & \nt{NORA}: la ganancia $g$ elige entre buscar y decidir (capítulo~\ref{sec:nora}) & Aumento de la relación señal/ruido: en el mono despierto, la noradrenalina en la corteza auditiva suprime la actividad de fondo más que la respuesta a las voces de congéneres (medido); la ganancia multiplicativa es un modelo; el papel en la elección entre buscar y decidir, una teoría & \cite{Foote1975NE,ServanSchreiber1990,AstonJones2005} & hipótesis \\
10 & \nt{ACET}: alterna escritura y lectura (capítulo~\ref{sec:acet}) & Las tres acciones (debilitar las conexiones internas, reforzar la entrada, facilitar la plasticidad) están medidas; el cambio de modo lo apoya indirectamente un experimento con escopolamina en humanos & \cite{HasselmoBower1992,Gil1997,Atri2004} & hipótesis \\
11 & La fórmula~\eqref{eq:machine} describe la máquina completa & Es un resumen de los modelos de los capítulos~\ref{sec:glutcomputer}--\ref{sec:acet}; cada parte se apoya en su capítulo; como un todo no se ha puesto a prueba experimentalmente & deducción en el capítulo & hipótesis \\
12 & Las perillas trabajan de forma coordinada sin un control común: error, sorpresa, escritura, regreso (fig.~\ref{fig:timeline}) & No hay experimento: el esquema es cualitativo. Eslabones sueltos: la teoría del reparto de tareas entre \nt{NORA} y \nt{ACET} y la relación entre la pupila y la tasa de aprendizaje en humanos & \cite{YuDayan2005,Nassar2012} & hipótesis \\
13 & El aprendizaje con una única señal común es tanto más lento cuantas más neuronas influyen en el resultado (proposición~\ref{prop:broadcastcost}) & Cálculo de las curvas de aprendizaje de una red lineal: el aprendizaje por perturbación de las salidas es más lento que el descenso por gradiente tantas veces como salidas haya & \cite{Werfel2005} & teoría \\
14 & No se sabe cómo la corteza ajusta los circuitos de muchas capas; candidatos: ráfagas de espigas, cálculos en las ramas de la neurona, autoaprendizaje por predicción & Las ráfagas como portadoras del error: modelo; solo una red de $5$--$8$ capas logró reproducir la respuesta de un modelo detallado de neurona cortical: modelo; espigas de calcio en las ramas de neuronas corticales humanas que dan el OR exclusivo: medido en cortes; autoaprendizaje: revisión de pruebas & \cite{Payeur2021,Beniaguev2021,Gidon2020,Keller2018} & hipótesis \\
\bottomrule
\end{longtable}}

\section{Capítulo~\ref{sec:sensors}. Las entradas de la máquina}

La estructura de los sensores se ha medido directamente: con registros de corriente de células individuales, de las vibraciones de la cóclea y con recuentos de células de la retina; el límite de sensibilidad y la fórmula del ruido de disparo se deducen de la física, y que los códigos de los sensores estén ajustados para la máxima información es una hipótesis con un apoyo sólido.

{\footnotesize
\setlength{\tabcolsep}{3pt}\renewcommand{\arraystretch}{1.15}
\begin{longtable}{>{\raggedright}p{0.5cm}>{\raggedright}p{4.0cm}>{\raggedright}p{5.6cm}>{\raggedright}p{2.3cm}>{\raggedright\arraybackslash}p{1.7cm}}
\toprule
N.º & Afirmación & Experimento y qué se midió & Fuente & Estado \\
\midrule\endfirsthead
\toprule
N.º & Afirmación & Experimento y qué se midió & Fuente & Estado \\
\midrule\endhead
1 & El sensor es un transductor: el receptor produce una corriente, y la salida de las células sensoriales del ojo y del oído es glutamato hacia el bus \nt{GLUT}; las primeras células no emiten espigas (fig.~\ref{fig:transducer}) & Los bastones y conos de la tortuga liberan un aminoácido excitador: lo detectan canales de glutamato en un parche de membrana escindido. Células ciliadas internas de la cóclea de rata: las corrientes en las terminales de las fibras nerviosas pasan por receptores AMPA de glutamato, y la frecuencia de liberación crece con la despolarización gradual de la célula hasta $150$\,Hz & \cite{CopenhagenJahr1989,GlowatzkiFuchs2002} & medido \\
2 & El sensor de luz en la oscuridad responde a un solo fotón con una corriente de alrededor de un picoamperio & Electrodo de succión sobre el segmento externo de un bastón de sapo: las respuestas a destellos tenues están cuantizadas, evento de $\approx1$\,pA ($5\%$ de la corriente oscura) con dispersión de $0.2$\,pA, eficiencia de $0.5$. Las personas informan de la llegada de un solo fotón más a menudo que al azar & \cite{Baylor1979,Tinsley2016} & medido \\
3 & El sensor de sonido detecta desplazamientos de menos de un nanómetro & Método Mössbauer, membrana basilar de la cóclea del cobayo en el punto de $16$--$19$\,kHz: en el umbral de respuesta del nervio auditivo, la velocidad de la membrana es de unos $0.04$\,mm/s, es decir, un desplazamiento de $\approx0.35$\,nm (cálculo nuestro). Se midió la membrana, no el haz del sensor & \cite{Sellick1982,Hudspeth1989} & medido \\
4 & En la oscuridad la precisión está limitada por el ruido de disparo de los fotones~\eqref{eq:photonnoise}; con luz débil la acumulación es más larga & La fórmula se deduce de la estadística de Poisson. En el bastón, el ruido con luz tenue coincide con la suma de eventos cuánticos independientes; sobre un fondo brillante, la respuesta a un destello es menor y más breve. La cifra de “décimas de segundo” no se comprobó aquí por separado & \cite{Baylor1979} & teoría \\
5 & El rango de entrada es de diez órdenes para la luz y doce para el sonido, y el de la frecuencia de espigas, de menos de tres órdenes & Para el sonido: la respuesta de la membrana basilar a la frecuencia característica crece con una pendiente de hasta $0.2$\,dB/dB en la banda de $40$--$80$\,dB, y la compresión se ve también por encima de $100$\,dB. Los números de órdenes son estimaciones estándar, sin fuente en el capítulo & \cite{Ruggero1997} & medido \\
6 & Respuesta del sensor~\eqref{eq:nakarushton}, con $\sigma$ siguiendo el brillo medio y desplazando la curva por el eje logarítmico (fig.~\ref{fig:adaptation}) & La fórmula se ajustó por primera vez a los potenciales S de la retina de pez. Conos de tortuga: las respuestas obedecen $V=I V_m/(I+\sigma)$ y abarcan $\sim3.5$ órdenes de brillo; tras $2$--$3$ minutos de fondo, la curva se desplaza horizontalmente sin cambiar de ancho. La relación con la división por la actividad total es una revisión & \cite{NakaRushton1966,NormannPerlman1979,Carandini2012} & medido \\
7 & Los sensores transmiten cambios, y sus códigos están ajustados para la máxima información por espiga (proposición~\ref{prop:sensors}) & El principio se propuso teóricamente. La prueba más fuerte: en la mosca, la característica “contraste $\to$ respuesta” de las neuronas de la primera capa coincide con la distribución acumulada de contrastes de las escenas naturales, y así se alcanza la máxima información & \cite{Barlow1961,Laughlin1981} & hipótesis \\
8 & Los sensores de encendido y apagado son un código de dos cables; la cámara de eventos lo repite en silicio & Revisión de experimentos: los canales de encendido y de apagado de la retina se separan ya en la primera sinapsis y pueden anularse por separado. Cámara: $128\times128$ píxeles sin fotogramas, rango de $120$\,dB, latencia de $15$\,\textmu s & \cite{Schiller1992,Lichtsteiner2008,Gallego2022} & medido \\
9 & En el ojo hay $\sim10^8$ sensores de luz y $\sim10^6$ neuronas de salida: compresión de $\sim100$ veces & Recuento en retinas humanas completas: en promedio $92$ millones de bastones y $4.6$ millones de conos; células de salida (ganglionares), de $0.7$ a $1.5$ millones según el ojo & \cite{Curcio1990,CurcioAllen1990} & medido \\
10 & En el ratón, el ojo tiene más de treinta canales de salida & Ratón: imagen de calcio con dos fotones de más de $11\,000$ células de salida y agrupamiento de sus respuestas: bastante más de $30$ tipos funcionales. En humanos el número no se ha medido & \cite{Baden2016} & medido \\
11 & El ojo del cobayo transmite $\sim875\,000$ bit/s; el cálculo para el ojo humano da $\sim10^7$ bit/s & Cobayo, matriz multielectrodo, movimiento en escenas naturales: $6$--$13$ bit/s por célula, $\sim875\,000$ bit/s para $\sim10^5$ células de salida. Para el ser humano ($\sim10^6$ células), $\sim10^7$ bit/s es un cálculo de los autores & \cite{Koch2006} & medido \\
12 & El oído es un banco de filtros pasa banda con un mapa de frecuencias casi logarítmico (fig.~\ref{fig:filterbank}) & El lugar de máxima vibración de la membrana basilar depende de la frecuencia; la función “frecuencia--lugar” es casi exponencial y describe con una sola forma las cócleas humanas y de animales vivos & \cite{Greenwood1990,Robles2001} & medido \\
13 & Los resonadores son activos: parte de los sensores amplifica las vibraciones débiles miles de veces en amplitud ($66$--$76$\,dB), y la ganancia cae al aumentar el volumen & Vibrometría láser de la base de la cóclea de la chinchilla: con tonos débiles a la frecuencia característica, ganancia de $66$--$76$\,dB respecto del estribo; la muerte reduce la sensibilidad en $60$--$81$\,dB ($10^3$--$10^4$ veces en amplitud) y elimina la compresión. La fuente de fuerza son las células ciliadas externas & \cite{Ruggero1997,Robles2001} & medido \\
14 & El amplificador del oído está al borde de la autooscilación & Cálculo: cerca de un punto de bifurcación de Hopf son inevitables la compresión del rango, la sintonía aguda y los tonos de combinación, todas las no linealidades conocidas de la audición. Que la cóclea esté ajustada así no se ha demostrado directamente & \cite{Eguiluz2000} & hipótesis \\
15 & A bajas frecuencias las espigas van en fase con el sonido (en el cobayo el enganche se debilita desde $\sim600$\,Hz y se nota hasta $\sim3.5$\,kHz), y con ellas se halla la dirección; a altas queda el código de lugar & Nervio auditivo del cobayo: el enganche de fase empieza a caer hacia $600$\,Hz y desaparece por encima de $3.5$\,kHz; en el gato y el mono el límite está aproximadamente una octava más arriba. La diferencia de tiempo de llegada a los dos oídos: revisión & \cite{PalmerRussell1986,Grothe2010} & medido \\
16 & Los demás sensores (tabla~\ref{tab:sensors}): el olfato tiene $\sim400$ tipos de receptores, uno por neurona, y un código combinatorio; el gusto funciona en parte con ATP y serotonina; el tacto, con sensores de adaptación rápida y lenta; calor y frío, por separado & Ratón: un receptor reconoce muchos olores, un olor activa muchos receptores, y olores distintos activan conjuntos distintos. Sin receptores de ATP el nervio del gusto no responde; los botones gustativos liberan serotonina. Registro de fibras individuales de la piel de la mano humana: cuatro tipos según la velocidad de adaptación. El canal del frío es distinto de los del calor & \cite{Mombaerts2004,Malnic1999,Finger2005,Huang2005,JohanssonVallbo1979,McKemy2002} & medido \\
\bottomrule
\end{longtable}}

\section{Capítulo~\ref{sec:recognition}. Reconocimiento}

La velocidad del reconocimiento, la estructura de las primeras etapas, la lectura lineal de la respuesta y el aprendizaje a partir de la continuidad del tiempo están medidos en humanos y en monos; la reducción de toda la cadena a dos operaciones y el aprendizaje a partir de video se comprobaron en los modelos del libro, y las explicaciones de las ilusiones van de bien fundadas a discutidas.

{\footnotesize
\setlength{\tabcolsep}{3pt}\renewcommand{\arraystretch}{1.15}
\begin{longtable}{>{\raggedright}p{0.5cm}>{\raggedright}p{4.0cm}>{\raggedright}p{5.6cm}>{\raggedright}p{2.3cm}>{\raggedright\arraybackslash}p{1.7cm}}
\toprule
N.º & Proposición & Experimento y qué se midió & Fuente & Estado \\
\midrule\endfirsthead
\toprule
N.º & Proposición & Experimento y qué se midió & Fuente & Estado \\
\midrule\endhead
1 & Con traslaciones, la comparación píxel a píxel con un patrón no acierta más que una moneda; el par de etapas~\eqref{eq:scstage} reconoce la figura en cualquier lugar & \texttt{models/\allowbreak recognition\_models.py}, “retina” de $24$ puntos, $4000$ pruebas: patrón, $0.49$, $0.51$, $0.52$ con fracciones de puntos invertidos $0$, $0.1$, $0.2$; par $S$ y $C$, $1.00$, $0.86$, $0.70$ & deducción en el capítulo & modelo \\
2 & Un animal en una imagen se reconoce en $\sim150$\,ms, una pasada por una decena de etapas de $10$--$20$\,ms & Humanos, tarea “hay un animal o no” con fotografías mostradas $20$\,ms: los potenciales evocados de las dos respuestas divergen a los $\sim150$\,ms. Mono: el tiempo de la primera respuesta crece de etapa en etapa (V1 antes que todas, luego V2 y V4). El número de etapas y “cero o una espiga por etapa” se deducen de estas cifras & \cite{Thorpe1996,Schmolesky1998} & medido \\
3 & La etapa $S$ es un AND de partes; la etapa $C$, el máximo sobre posiciones (proposición~\ref{prop:conveyor}) & Gato, corteza visual primaria: las células simples responden a un borde de cierta orientación en cierto lugar; las complejas, a la misma orientación en una zona mayor. Registro intracelular de células complejas: en muchas, la respuesta a dos barras se acerca a la mayor de las dos respuestas, en promedio lo más cerca de la operación max. El máximo mediante inhibición mutua es la proposición~\ref{prop:wta}; la división por la actividad total, una revisión & \cite{HubelWiesel1962,Lampl2004,Carandini2012} & medido \\
4 & Más arriba en la cadena, las combinaciones son más grandes y complejas, y la respuesta depende cada vez menos de la posición & Mono, simplificación de estímulos hasta el “rasgo crítico”: las células de V4 y de la corteza temporal inferior posterior responden a combinaciones complejas de forma, color y textura con campo pequeño; en la corteza temporal inferior anterior el campo es mayor. La alternancia de $S$ y $C$ en un modelo reproduce este cuadro & \cite{KobatakeTanaka1994,Riesenhuber1999} & medido \\
5 & Las redes convolucionales repiten esta alternancia y predicen mejor que nada las respuestas de las etapas superiores; el objeto se lee linealmente de las frecuencias de un grupo de neuronas & Una red entrenada para reconocer, sin ajuste al cerebro: la capa superior predice las respuestas de las neuronas de la corteza temporal inferior del mono; las intermedias, las respuestas de V4. Un clasificador lineal sobre la actividad de $\sim100$ células elegidas al azar, en ventanas desde $12.5$\,ms, reconoce el objeto y la categoría con distintas posiciones y tamaños & \cite{Fukushima1980,Yamins2014,Hung2005} & medido \\
6 & La regla de Hebb con traza~\eqref{eq:traceinvariance} enseña la invariancia a partir de video sin etiquetas si la traza no es más corta que $\sim20$\,ms & La idea de los rasgos lentos y de la regla con traza es teoría. \texttt{models/\allowbreak recognition\_models.py}, dos neuronas $C$, $16$ entradas, $10$ corridas, pausa entre objetos de $0.3$\,s. Con $\tau_{\text{tr}}=5$\,ms la coincidencia con la figura es de $0.52$ en promedio, con $10$\,ms, $0.56$; con $20$, $50$ y $200$\,ms, $1.00$ en todas las corridas (con $30$\,ms, una corrida de diez se equivoca) & \cite{Foldiak1991,WiskottSejnowski2002} & modelo \\
7 & En el cerebro la invariancia se aprende de la continuidad del tiempo & Mono: el objeto se sustituía sin que se notara durante un salto del ojo hacia un lugar determinado del campo; la invariancia de posición de las neuronas de la corteza temporal inferior cambiaba de acuerdo con la sustitución, de forma significativa ya al cabo de una hora. Que sea la regla principal del aprendizaje visual es una hipótesis & \cite{LiDiCarlo2008} & medido \\
8 & Movimiento: detectores de dirección y luego el mismo par AND/OR sobre zonas grandes & Detectores de dirección en la retina del conejo; en el área MT del mono, las $110$ neuronas registradas son selectivas a la dirección, y la respuesta al movimiento es más fuerte que en V1 & \cite{Barlow1965,Albright1984} & medido \\
9 & Las etapas superiores envían hacia abajo una predicción, y hacia arriba va el error & El modelo explica los efectos extraclásicos de los campos de la corteza visual primaria; algunas observaciones concuerdan (revisión); como estructura general de la cadena, no está demostrado & \cite{RaoBallard1999,Keller2018} & hipótesis \\
10 & Las imágenes difíciles requieren realimentación y varias vueltas & Mono: para las imágenes “difíciles”, la decisión en la corteza temporal inferior se forma $\sim30$\,ms más tarde; estas respuestas tardías las predice mal una red directa y mejor las redes con realimentación o más profundas & \cite{Kar2019} & medido \\
11 & Una alteración de píxeles invisible engaña a la red; la visión humana es mucho más robusta, pero no invulnerable & En redes: una perturbación pequeña, elegida a propósito, cambia la respuesta; el ejemplo “panda $\to$ gibón” es del segundo trabajo. En humanos, con presentación breve, las falsificaciones que se transfieren bien entre redes desvían la elección & \cite{Szegedy2014,Goodfellow2015,Elsayed2018} & medido \\
12 & Ilusiones de expectativa: la entrada poco fiable cede ante la expectativa; lo pálido se mueve “más despacio”, el vestido se ve de formas distintas (sección~\ref{sec:illusions}) & Las rejillas de menor contraste parecen moverse más despacio. Encuesta a $1401$ personas: el vestido se ve en tres variantes estables, y una pista sobre la iluminación cambia el color; en $\sim13\,000$ observadores decide la suposición sobre la iluminación. Un modelo bayesiano con expectativa de movimientos lentos explica varias ilusiones de movimiento & \cite{Thompson1982,LaferSousa2015,Wallisch2017,Weiss2002,Purves2011} & hipótesis \\
13 & Control de ganancia: la cascada, la inclinación y el contraste; la rejilla de Hermann no es inhibición de vecinos & Los detectores de dirección de la retina del conejo reducen su frecuencia ante un movimiento prolongado y, al detenerse, caen por debajo del fondo. La ilusión de inclinación se reproduce con un modelo de división por la actividad de los vecinos. Modificaciones de la rejilla que no cambian la respuesta de las neuronas de salida de la retina eliminan las manchas: se descarta la explicación por inhibición de vecinos & \cite{BarlowHill1963,Schwartz2009,SchillerCarvey2005} & medido \\
14 & Compensación de retardos: un objeto en movimiento parece ir por delante del destello & La ilusión está medida. La psicofísica contradice tanto la prolongación del movimiento hacia adelante como la diferencia de latencias; se propuso una explicación retrospectiva, según los sucesos de $\sim80$\,ms después del destello. El debate no está resuelto & \cite{Nijhawan1994,EaglemanSejnowski2000} & hipótesis \\
15 & Las “serpientes” inmóviles giran: los detectores de dirección leen la diferencia de latencias~\eqref{eq:latency} & La ilusión funciona también sin color; los pares de elementos vecinos la generan por separado; las neuronas corticales del mono selectivas a la dirección dan una respuesta direccional a estos pares inmóviles. En humanos, el giro lo disparan las microsacadas y los parpadeos & \cite{Conway2005,OteroMillan2012} & medido \\
16 & Imágenes ambiguas: inhibición mutua, fatiga y ruido; los cambios los provoca sobre todo el ruido & Modelo de grupos de neuronas en competencia: en él los cambios los provoca el \emph{ruido}, y sin ruido no ocurren; explica el aumento de la frecuencia de cambios con la fuerza del estímulo. La fatiga tiene aquí un papel menor & \cite{MorenoBote2007} & hipótesis \\
17 & Parte de las ilusiones es consecuencia de la tarea que aprende la máquina & Una red entrenada para predecir el cuadro siguiente de un video “ve” el giro de los anillos inmóviles y no lo ve en las imágenes de control; las redes para limpiar imágenes y aumentar su nitidez caen en ilusiones de color y brillo, como las personas & \cite{Watanabe2018,GomezVilla2020} & medido \\
\bottomrule
\end{longtable}}

\section{Capítulo~\ref{sec:muscles}. La salida a los músculos}

La estructura de la salida (unidades motoras, margen de seguridad de la sinapsis, activación por tamaño, par de músculos, lazo de estiramiento y generadores de ritmo) está medida directamente; la condición de estabilidad del lazo con retardo es un teorema; el modelo interno del cuerpo es una hipótesis bien fundada; los números de las figuras son cálculos con el modelo del libro.

{\footnotesize
\setlength{\tabcolsep}{3pt}\renewcommand{\arraystretch}{1.15}
\begin{longtable}{>{\raggedright}p{0.5cm}>{\raggedright}p{4.0cm}>{\raggedright}p{5.6cm}>{\raggedright}p{2.3cm}>{\raggedright\arraybackslash}p{1.7cm}}
\toprule
N.º & Afirmación & Experimento y qué se midió & Fuente & Estado \\
\midrule\endfirsthead
\toprule
N.º & Afirmación & Experimento y qué se midió & Fuente & Estado \\
\midrule\endhead
1 & Unidad motora: una neurona \nt{ACET} controla un grupo de fibras, de unos cientos a un par de miles; en los músculos de ajuste fino las unidades son más pequeñas & Músculos humanos, recuento de axones motores y de fibras musculares: en promedio unas $340$ fibras por neurona en un músculo interóseo de la mano, unas $410$ en el braquiorradial, unas $1930$ en el gastrocnemio medial. El capítulo no da una cifra exacta para las unidades más pequeñas. & \cite{Feinstein1955mu} & medido \\
2 & La sinapsis sobre el músculo libera varias veces más neurotransmisor del necesario para activar la fibra & Revisión de mediciones en sinapsis neuromusculares: el margen de seguridad (cociente entre el neurotransmisor liberado y el umbral) en mamíferos adultos suele ser de $3$--$5$; en humanos, el margen se apoya más en los pliegues de la membrana receptora que en la cantidad liberada. & \cite{WoodSlater2001} & medido \\
3 & La sacudida~\eqref{eq:twitch} sube en un tiempo $T_c$ del orden de $50\ms$ & Unidades motoras individuales de la mano humana con esfuerzo voluntario, fuerza promediada sobre las espigas de la unidad: tiempo de subida de $30$--$100\ms$, menos de $70\ms$ en el $80\,\%$ de las unidades. La función $(t/T_c)e^{1-t/T_c}$ es una aproximación tomada de un modelo del conjunto de unidades. & \cite{MilnerBrown1973rec, Fuglevand1993} & medido \\
4 & Suma de sacudidas~\eqref{eq:forcesum}: con $5$, $15$ y $40$ espigas/s, fuerza de $0.2$--$1.1F_1$, $1.8$--$2.2F_1$ y cerca de $5.4F_1$ (fig.~\ref{fig:tetanus}) & Cálculo con \texttt{models/\allowbreak muscle\_\allowbreak models.py}, $T_c=50\ms$: $0.21$--$1.10$, $1.76$--$2.19$ y $5.28$--$5.49\,F_1$ en los últimos $0.3\,\text{s}$. La suma lineal es una simplificación: el modelo no incluye la saturación del músculo real. & \texttt{models/\allowbreak muscle\_\allowbreak models.py} & modelo \\
5 & Las unidades se activan por tamaño; las más débiles son cien veces más débiles que las más fuertes & Raíces ventrales del gato: con una entrada refleja creciente, las primeras en descargar son las neuronas con espigas pequeñas en la raíz, es decir, las pequeñas. Músculo interóseo de la mano humana: fuerza de sacudida de las unidades de $0.1$ a $10$~g, que crece casi linealmente con la fuerza a la que se activa la unidad. & \cite{Henneman1957, MilnerBrown1973rec} & medido \\
6 & El paso de fuerza es proporcional a la fuerza, $\Delta F/F\approx q-1$~\eqref{eq:sizeprinciple}: punto flotante & Deducido en el capítulo a partir de la progresión geométrica de las fuerzas. Concuerda con los experimentos: con esfuerzos grandes, para el mismo incremento de fuerza se activan muchas menos unidades nuevas. & deducción en el capítulo; \cite{MilnerBrown1973rec} & teoría \\
7 & La dispersión de la fuerza crece en proporción a la propia fuerza & Esfuerzo isométrico voluntario de un músculo del pulgar: la desviación estándar de la fuerza crece linealmente con la fuerza media; con el mismo esfuerzo provocado por estimulación eléctrica, la dispersión es constante. Modelo del conjunto de unidades: el crecimiento lineal lo produce la activación de las unidades por orden de fuerza. & \cite{Jones2002sdn} & medido \\
8 & La suavidad de los movimientos es consecuencia del ruido dependiente de la señal & Modelo: la trayectoria que minimiza la dispersión del punto final con ese ruido reproduce las trayectorias medidas del ojo y del brazo. No hay un experimento directo que distinga esta causa de otras. & \cite{HarrisWolpert1998} & hipótesis \\
9 & La diferencia de fuerzas del par de músculos fija el momento; la suma, la rigidez~\eqref{eq:antagonists} & Teoría del control de la impedancia mediante cocontracción. Brazo humano: la rigidez de cada articulación, medida con pequeños empujones, es aproximadamente proporcional a su momento; distintas proporciones de cocontracción cambian la forma y la dirección de la elipse de rigidez de la mano. & \cite{Hogan1984, GomiOsu1998stiff} & medido \\
10 & El sensor de estiramiento excita su propia neurona \nt{ACET} y, a través de una intermediaria inhibidora, apaga al antagonista (fig.~\ref{fig:antagonists}) & Médula espinal del gato: una sola intermediaria, excitada por las fibras de los sensores de estiramiento, produce potenciales inhibidores monosinápticos en las neuronas \nt{ACET} del músculo antagonista, con un retardo sináptico de $0.28$--$0.42\ms$. El neurotransmisor de la intermediaria (glicina) no se determinó en este experimento. & \cite{Jankowska1972Ia} & medido \\
11 & Retardo del lazo: espinal, $25$--$30\ms$; a través de la corteza, entre una vez y media y tres veces más; a través del ojo, $100$--$200\ms$ & Brazo humano, empujón a la articulación: respuesta muscular corta a los $20$--$45\ms$, larga a los $45$--$105\ms$, voluntaria a los $120$--$180\ms$. Desplazamiento de la posición visible del dedo durante el movimiento: corrección a los $160\ms$ en promedio. & \cite{Pruszynski2008rapid, SaundersKnill2003vis} & medido \\
12 & $\dd x/\dd t=-Gx(t-d)$ es estable con $Gd<\pi/2$~\eqref{eq:delaystable}; en el límite la frecuencia es $1/(4d)$ & Teorema sobre las raíces de una ecuación con retardo. Comprobación con \texttt{models/\allowbreak muscle\_\allowbreak models.py}, $d=30\ms$: a los $3\,\text{s}$ la amplitud es $3\cdot10^{-6}$ con $G=40$, $0.13$ con $G=50$ y $38$ con $G=55$ por segundo (límite: $52$). & \cite{Hayes1950} & teoría \\
13 & Temblor fisiológico de $8$--$12$~Hz; el lazo de estiramiento es una de sus fuentes & Revisión de mediciones: las personas sanas tienen un temblor fino en torno a $10$~Hz; su origen es mixto (oscilaciones centrales, propiedades de las descargas de las unidades, resonancia mecánica y resonancia del lazo reflejo), y en distintas condiciones predomina uno u otro. & \cite{McAuleyMarsden2000} & hipótesis \\
14 & El cerebro predice el resultado de una orden con un modelo interno del cuerpo & Una persona sostiene una carga en pinza y mueve el brazo: con carga inercial, viscosa y elástica, la fuerza de prensión cambia al mismo tiempo que la carga, no con el retardo del lazo, es decir, la carga se predice. Una revisión reúne estos experimentos; no hay observación directa del modelo mismo en las neuronas. & \cite{FlanaganWing1997grip, WolpertGhahramani2000} & hipótesis \\
15 & El ritmo de caminar, respirar y masticar lo generan redes locales con una entrada constante & Revisión de experimentos: un sistema nervioso aislado (médula espinal, ganglios) produce un patrón motor rítmico sin entrada rítmica y sin realimentación de los músculos. & \cite{MarderBucher2001} & medido \\
16 & La inhibición mutua con fatiga~\eqref{eq:halfcentre} da un ritmo de período $0.84\,\text{s}\approx2\tau_a$ (fig.~\ref{fig:halfcentre}) & Teorema: una red con inhibición mutua y adaptación sin estado estable oscila necesariamente con entrada constante. Cálculo con \texttt{models/\allowbreak muscle\_\allowbreak models.py} ($I=1$, $\beta=2$, $b=2$, $\tau=20\ms$, $\tau_a=0.4\,\text{s}$): período de $0.84\,\text{s}$. No se ha demostrado que los generadores reales estén construidos así. & \cite{Matsuoka1985osc} & modelo \\
\bottomrule
\end{longtable}}

\section{Capítulo~\ref{sec:pain}. El dolor}

Los sensores del dolor, los dos cables, la compuerta en la médula espinal, la perilla descendente de volumen, la sensibilización y la captura de la atención están medidos directamente; las fórmulas de la compuerta y de la sensibilización son modelos simplificados del libro; la regla con la que la máquina elige el volumen de la alarma es una hipótesis.

{\footnotesize
\setlength{\tabcolsep}{3pt}\renewcommand{\arraystretch}{1.15}
\begin{longtable}{>{\raggedright}p{0.5cm}>{\raggedright}p{4.0cm}>{\raggedright}p{5.6cm}>{\raggedright}p{2.3cm}>{\raggedright\arraybackslash}p{1.7cm}}
\toprule
N.º & Proposición & Experimento y qué se midió & Fuente & Estado \\
\midrule\endfirsthead
\toprule
N.º & Proposición & Experimento y qué se midió & Fuente & Estado \\
\midrule\endhead
1 & Umbral alto: los umbrales de calor de los distintos sensores del dolor van de $41^\circ$ a $46$--$48^\circ$ & Registro de fibras individuales. Piel de mono: umbral mediano de los sensores lentos (C), $41^\circ$; de los sensores rápidos de tipo~II, $46^\circ$; de tipo~I, más de $53^\circ$. Piel humana: sensores C que responden también a la presión, $40.7^\circ$; los que no responden a la presión, $48^\circ$. & \cite{Treede1995heat, Weidner1999cfib} & medido \\
2 & Los sensores del dolor se adaptan mucho menos que los del tacto y tras una lesión leve se vuelven más sensibles; el debilitamiento tras un calentamiento intenso se midió en un tipo & Quemadura leve de la piel ($50^\circ$, $100\,\text{s}$): a los $5$--$10$~min el umbral del dolor en personas y el umbral de los sensores C en el mono bajan $2$--$6^\circ$; otros sensores de la piel no muestran ese desplazamiento. En los sensores rápidos de tipo~II la respuesta tras un calentamiento intenso queda suprimida. La comparación de la adaptación con los sensores del tacto es una estimación general, sin un experimento propio que la respalde aquí. & \cite{LaMotte1982hyper, Treede1995heat} & medido \\
3 & Dos cables: rápido, $5$--$30$~m/s; lento, $0.5$--$2$~m/s; tiempo de llegada $t=L/v$~\eqref{eq:paintime} & Velocidad de conducción: sensores rápidos del dolor del mono, en promedio $15$ y $25$~m/s (dos tipos); sensores C humanos, $0.8$--$1.0$~m/s. Con $L=1$~m son decenas de milisegundos y alrededor de un segundo. & \cite{Treede1995heat, Weidner1999cfib} & medido \\
4 & El dolor llega dos veces: primero rápido y preciso, luego sordo y prolongado & Tiempo de reacción al calentamiento del antebrazo humano: de $1100$ a $400\ms$ al subir la temperatura; un calentamiento intenso de la piel con vello produce doble dolor, y el de la palma, no. El primer dolor solo pueden llevarlo fibras de más de $6$~m/s; por su latencia corresponde a los sensores rápidos de tipo~II, que la palma no tiene. El reparto de papeles (“dónde” y “qué tan grave”) es una interpretación. & \cite{CampbellLaMotte1983first, Treede1995heat} & medido \\
5 & Compuerta: la neurona de salida de la médula espinal recibe dolor y tacto, y el tacto excita a la intermediaria inhibidora (fig.~\ref{fig:gate}) & El esquema de la compuerta se propuso como teoría. Ratones, marcaje genético y eliminación de grupos de neuronas: las neuronas excitadoras con somatostatina son necesarias para el dolor mecánico; las neuronas inhibidoras con dinorfina impiden que la entrada de los sensores del tacto llegue a ellas, y sin estas neuronas un roce leve provoca dolor. & \cite{MelzackWall1965, Duan2014} & medido \\
6 & El tacto atenúa el dolor & Personas, pulsos láser dolorosos en la mano: un tacto simultáneo reduce la valoración del dolor y la capacidad de distinguir la energía de los pulsos; el efecto se debilita con la distancia entre el punto del tacto y el del dolor dentro de un mismo segmento cutáneo. & \cite{Mancini2014touch} & medido \\
7 & Supuesto de la teoría de la compuerta: el dolor fuerte apaga por sí mismo a la intermediaria inhibidora, y la compuerta se abre de par en par & En el capítulo se señala como supuesto de la teoría original. El trabajo en ratones describe cómo la compuerta impide que el tacto entre en el canal del dolor; no muestra que el dolor apague a la intermediaria. & \cite{MelzackWall1965} & hipótesis \\
8 & Compuerta~\eqref{eq:gate}: umbral de $0.18$ sin tacto, $0.86$ con tacto, $0.11$ con $g=2$; con $\nu=1$ el tacto reduce la salida de $1$ a $0.25$, con $\nu=2$ no actúa; el tacto solo no pasa (fig.~\ref{fig:gatecurves}) & El cálculo con \texttt{models/\allowbreak pain\_\allowbreak gate.py} con los pesos del texto reproduce todos estos números. & \texttt{models/\allowbreak pain\_\allowbreak gate.py} & modelo \\
9 & Las vías descendentes del tronco cambian la ganancia de la compuerta en ambos sentidos & Rata, bulbo raquídeo, registro durante el calentamiento de la cola: unas neuronas descargan antes de la retirada (“on”), otras se callan (“off”); una estimulación débil de esta zona suprime el reflejo. Revisión: los mismos circuitos facilitan e inhiben la transmisión del dolor, y los opioides actúan a través de ellos. & \cite{Fields1983rvm, Fields2004} & medido \\
10 & La expectativa de alivio atenúa el dolor a través del canal opioide & Personas con dolor agudo: en aquellas a quienes la expectativa de alivio les redujo el dolor, el bloqueo de los receptores opioides devolvió el dolor al nivel de las demás; en las demás, el bloqueo no cambió el dolor. & \cite{Levine1978} & medido \\
11 & El cerebro elige el volumen de la alarma según el beneficio de la interrupción & No hay experimento que haya medido la regla de elección del volumen. & --- & hipótesis \\
12 & Sensibilización: tras una lesión, la salida de la compuerta crece y el umbral baja, en parte por la propia médula espinal; persiste después de que cesa el estímulo lesivo & Rata: tras una lesión de la piel, el umbral del reflejo flexor baja y la respuesta crece; el análisis electrofisiológico muestra que parte del aumento surge en la propia médula espinal. El estímulo lesivo provoca alteraciones duraderas de la sensibilidad que sobreviven al propio estímulo. El capítulo no da el tiempo exacto de recuperación. & \cite{Woolf1983} & medido \\
13 & La sensibilización~\eqref{eq:sensitization} es una realimentación positiva; con un lazo fuerte, la alarma puede trabarse después de la curación & Se deduce del modelo~\eqref{eq:sensitization} (ejercicio al final del capítulo). No hay experimento que muestre que el dolor persistente tras la curación sea un lazo trabado precisamente de este tipo. & deducción en el capítulo & hipótesis \\
14 & El dolor es una recompensa negativa, y el aprendizaje con él sigue el error de diferencias temporales & Resonancia, personas, cadenas de precursores del dolor: la actividad del estriado ventral y de la ínsula anterior coincide con las señales de aprendizaje por diferencias temporales. Mono: parte de las neuronas \nt{DOPA} se excita con un soplo de aire desagradable y su precursor; otra parte se inhibe. & \cite{Seymour2004, Matsumoto2009} & medido \\
15 & El dolor interrumpe el trabajo en curso & Personas, estímulo eléctrico doloroso y sondas sonoras: la respuesta a una sonda $100\ms$ después del inicio del dolor es notablemente más lenta; a los $1.5\,\text{s}$ ya no hay diferencia con el control. Una revisión reúne estos experimentos en un modelo de captura de la atención. & \cite{Crombez1996disrupt, EcclestonCrombez1999} & medido \\
16 & La retirada se cierra en la médula espinal & Rata: las neuronas de las capas profundas del asta dorsal reproducen el patrón espacial del reflejo de retirada de cada músculo; no se logra excitarlas antidrómicamente desde el segmento cervical, es decir, son intermediarias espinales. & \cite{Schouenborg1995wr} & medido \\
\bottomrule
\end{longtable}}

\section{Capítulo~\ref{sec:motorlearning}. Aprendizaje motor}

La regla de mínimos cuadrados y la ventaja de un cable de error propio son teoremas; la estructura del circuito con cable de error, el debilitamiento por coincidencia, el ajuste de la traza al retardo, el efecto posterior y la recuperación espontánea de la habilidad están medidos; que todo el circuito funcione como aprendizaje supervisado y construya un modelo interno es la hipótesis principal; los números del modelo de dos procesos son cálculos con el modelo del libro.

{\footnotesize
\setlength{\tabcolsep}{3pt}\renewcommand{\arraystretch}{1.15}
\begin{longtable}{>{\raggedright}p{0.5cm}>{\raggedright}p{4.0cm}>{\raggedright}p{5.6cm}>{\raggedright}p{2.3cm}>{\raggedright\arraybackslash}p{1.7cm}}
\toprule
N.º & Afirmación & Experimento y qué se midió & Fuente & Estado \\
\midrule\endfirsthead
\toprule
N.º & Afirmación & Experimento y qué se midió & Fuente & Estado \\
\midrule\endhead
1 & La regla~\eqref{eq:lms} sigue en promedio el gradiente del cuadrado del error & Resultado de la teoría de los filtros adaptativos: una corrección proporcional a la entrada y al error de salida es un descenso estocástico por el gradiente del error cuadrático medio. & \cite{WidrowHoff1960} & teoría \\
2 & Con un cable de error por salida, la velocidad no depende del número de salidas; con un solo número común, cae con el número de neuronas ruidosas (proposición~\ref{prop:teachers}) & Curvas de aprendizaje de una red lineal: el aprendizaje por perturbaciones aleatorias de las neuronas es más lento que el gradiente exacto tantas veces como neuronas perturbadas haya. & \cite{Werfel2005} & teoría \\
3 & El circuito con cable de error funciona como aprendizaje supervisado (proposición~\ref{prop:errorwire}) & Teorías deducidas de la estructura del circuito. Los experimentos de las filas 7--10 concuerdan con ellas; que todo el circuito funcione exactamente así no está demostrado. & \cite{Marr1969, Albus1971} & hipótesis \\
4 & Capa de expansión: unos $7\cdot10^{10}$ pequeñas neuronas \nt{GLUT}, más que en todo el resto del cerebro & Recuento de núcleos en una suspensión de tejido disuelto: en el cerebelo humano hay $69\cdot10^9$ neuronas de las $86\cdot10^9$ de todo el cerebro. Estereología: $101\cdot10^9$ neuronas pequeñas en el cerebelo de hombres de edad avanzada. & \cite{Azevedo2009, Andersen1992cereb} & medido \\
5 & Elevar la dimensión de la entrada vuelve lineal la dependencia deseada & Teorema sobre el número de particiones: una suma ponderada de $n$ entradas con umbral realiza un etiquetado arbitrario de unos $2n$ vectores en posición general. Que el cerebelo aproveche justamente esto es parte de la hipótesis de la fila~3. & \cite{Cover1965} & teoría \\
6 & Unas $3\cdot10^7$ neuronas \nt{GABA} de salida, cada una con del orden de $10^5$ entradas & Estereología del cerebelo humano: $30.5\cdot10^6$ neuronas de salida. Microscopía electrónica, rata: unas $1.75\cdot10^5$ entradas de la capa de expansión por neurona de salida. El neurotransmisor inhibidor de las neuronas de salida, en la revisión. & \cite{Andersen1992cereb, NapperHarvey1988, Ito2001} & medido \\
7 & Exactamente un cable de error por neurona de salida, alrededor de una vez por segundo, cada vez una descarga potente & Revisión de registros: cada neurona de salida recibe una única entrada trepadora \nt{GLUT}, que provoca una espiga compleja con una frecuencia del orden de $1$~Hz. & \cite{Ito2001} & medido \\
8 & La probabilidad de la descarga depende de la dirección del error del movimiento & Mono, región oculomotora del cerebelo, adaptación de sacadas: la dirección del error visual fija la probabilidad de la espiga compleja, y su magnitud, su momento; la descarga cambia las espigas simples en el intento siguiente y desplaza el movimiento a lo largo del error preferido de la célula. & \cite{Herzfeld2018} & medido \\
9 & La coincidencia de la descarga de error con una entrada activa debilita la entrada ($-\eta e_i x_j$) & Revisión de experimentos: la estimulación conjunta de una entrada de la capa de expansión y del cable de error provoca un debilitamiento duradero de esa entrada; la estimulación de la entrada sola, no. & \cite{Ito2001} & medido \\
10 & La duración de la traza está ajustada al retardo del error: con un retardo de $100\ms$ se debilita sobre todo la entrada activa $100\ms$ antes & Cortes y ratones vivos: en la región responsable del aprendizaje oculomotor, el debilitamiento está finamente ajustado a un intervalo de unos $120\ms$ entre la entrada y el error, lo que tarda la señal de error en esa tarea; en otra región, distintas células se ajustan a distintos intervalos. & \cite{Suvrathan2016} & medido \\
11 & El circuito es un filtro adaptativo~\eqref{eq:adaptivefilter} que aprende un modelo interno del cuerpo & Modelo deducido de la estructura del circuito. No hay un experimento directo en el que el filtro aprendido se haya medido como función de transferencia del cuerpo. & \cite{Fujita1982} & hipótesis \\
12 & La predicción resta las consecuencias de los propios movimientos; por eso uno no puede hacerse cosquillas & Personas: el tacto propio resulta menos cosquilloso que el mismo tacto ajeno; en el escáner, la corteza somatosensorial responde más débilmente al tacto propio, y el cerebelo está menos activo cuando el movimiento toca la piel. La explicación por resta de la predicción es una hipótesis. & \cite{Blakemore1998} & hipótesis \\
13 & Con una fuerza externa el fallo disminuye, y al retirarla la mano falla hacia el otro lado & Personas que mueven el mango de un robot en un campo de fuerzas: las trayectorias vuelven a ser rectas; al retirar el campo, el efecto posterior es casi la imagen especular de las primeras distorsiones, y se transfiere también a zonas del espacio no entrenadas. & \cite{ShadmehrMussaIvaldi1994} & medido \\
14 & Aprenden dos procesos~\eqref{eq:tworate}; tras la perturbación inversa, la habilidad vuelve sola & Personas, campo de fuerzas, luego un breve campo inverso e intentos con el error fijado: la corrección vuelve sola a la anterior; el modelo de dos procesos lo reproduce, el de uno, no. & \cite{Smith2006} & medido \\
15 & Números de la fig.~\ref{fig:tworate}: $0.84$ (lento, $0.63$) en $200$ movimientos; $6$ inversos; rápido, $-0.53$; lento, $0.46$; regreso hasta $0.37$ en $40$ movimientos & Cálculo con \texttt{models/\allowbreak motor\_\allowbreak learning.py}, $A_f=0.92$, $B_f=0.1$, $A_s=0.996$, $B_s=0.02$: $0.840$ ($0.634$ y $0.205$), $6$ movimientos, $-0.531$ y $0.457$, pico de $0.371$ en el movimiento $40$. & \texttt{models/\allowbreak motor\_\allowbreak learning.py} & modelo \\
16 & Hacen falta dos procesos porque el cuerpo cambia a distintas velocidades & Teoría: la estimación óptima con perturbaciones de varias duraciones da varios filtros con distintas constantes de tiempo y explica la recuperación espontánea y el reaprendizaje acelerado. & \cite{Kording2007} & hipótesis \\
\bottomrule
\end{longtable}}

\section{Capítulo~\ref{sec:chemoutput}. La salida química}

Los dos cables hacia el corazón y su distinta velocidad, la realimentación negativa en las cascadas hormonales, el código por frecuencia de porciones, el generador diario y su ajuste por la luz están medidos directamente; el límite $K<8$ es un teorema de la teoría de control deducido en el capítulo; las pulsaciones generadas por la propia realimentación de la cascada son una hipótesis con un experimento sólido a su favor.

{\footnotesize
\setlength{\tabcolsep}{3pt}\renewcommand{\arraystretch}{1.15}
\begin{longtable}{>{\raggedright}p{0.5cm}>{\raggedright}p{4.0cm}>{\raggedright}p{5.6cm}>{\raggedright}p{2.3cm}>{\raggedright\arraybackslash}p{1.7cm}}
\toprule
N.º & Afirmación & Experimento y qué se midió & Fuente & Estado \\
\midrule\endfirsthead
\toprule
N.º & Afirmación & Experimento y qué se midió & Fuente & Estado \\
\midrule\endhead
1 & El marcapasos del corazón late por sí solo unas $100$ veces por minuto; en reposo el freno está pisado y la frecuencia suele ser del orden de $60$--$70$ & Humanos, bloqueo completo de ambos cables hacia el corazón: la frecuencia propia es mayor que la de reposo y disminuye con la edad; en adultos es del orden de $100$ latidos por minuto. Por tanto, en reposo predomina el freno. La frecuencia de reposo de $60$--$70$ es un valor típico, sin cita aparte. & \cite{JoseCollison1970ihr} & medido \\
2 & El acelerador \nt{NORA} y el freno \nt{ACET} son filtros pasa bajos, y el freno es más rápido~\eqref{eq:gasbrake} & Perro, estimulación modulada en frecuencia del nervio vago o simpático cardiaco y función de transferencia hacia la frecuencia auricular: ambas vías son filtros pasa bajos; la simpática tiene una frecuencia de corte mucho menor y un retardo puro añadido de unos $1.7\,\text{s}$. Las características dependen del tono medio, así que la fórmula lineal~\eqref{eq:gasbrake} es una aproximación. & \cite{Berger1989} & medido \\
3 & El freno es rápido porque \nt{ACET} abre el canal de potasio sin segundo mensajero, y \nt{NORA} lo hace mediante una cadena de reacciones & Parches de membrana de células auriculares: el canal de potasio que abre la acetilcolina lo abren las subunidades $\beta\gamma$ de la proteína G unida al receptor, sin segundo mensajero dentro de la célula. La vía de \nt{NORA} a través del AMPc no se cita aquí. & \cite{Logothetis1987gbg} & medido \\
4 & La sangre recorre el cuerpo en un tiempo del orden de un minuto; \nt{ADRE} llega desde el torrente sanguíneo a todos los órganos con receptor & Estimación general de orden de magnitud; así se formula en el capítulo, sin fuente citada. & --- & estimación \\
5 & La cascada hormonal está envuelta en una realimentación negativa: la hormona final inhibe las primeras etapas & Revisión de experimentos: las hormonas de la corteza suprarrenal suprimen la liberación de ACTH por la hipófisis y de la hormona liberadora por el hipotálamo en tres escalas de tiempo: rápida, intermedia y lenta. & \cite{KellerWood1984fb} & medido \\
6 & Tres etapas iguales son estables si $K<8$~\eqref{eq:cascadestable}; en la frontera el periodo es $2\pi\tau/\sqrt3\approx3.6\tau$ & Deducido en el capítulo: a la frecuencia $\sqrt3/\tau$ las tres etapas dan un desfase de $180^\circ$ y una atenuación de $8$ veces. & deducción del capítulo & teoría \\
7 & El error de nivel es $1/(1+K)$, por lo que un regulador así no mantiene el nivel con precisión mejor que un $10\%$ aproximadamente (proposición~\ref{prop:cascade}) & Consecuencia de la fila~6 para realimentación proporcional. El eje real tiene realimentación en varias etapas y en varias escalas de tiempo (fila~5), así que el límite se refiere al modelo~\eqref{eq:cascade} y no está medido. & deducción del capítulo & teoría \\
8 & Con $K=4$ y $7$ el nivel se estabiliza; con $K=12$, pulsaciones regulares con periodo de $3.7$--$3.9\,\tau$ (fig.~\ref{fig:cascade}) & Cálculo con \texttt{models/\allowbreak chem\_\allowbreak models.py}: con $K=12$ los periodos sucesivos son $3.9$, $3.7$, $3.8$, $3.7\,\tau$; con $K=4$ la amplitud al final del cálculo es $0.003$. & \texttt{models/\allowbreak chem\_\allowbreak models.py} & modelo \\
9 & La hormona del estrés se libera en pulsos más o menos una vez por hora; los pulsos los genera la propia realimentación, sin generador aparte & Ratas, infusión constante de hormona liberadora: la ACTH y la corticosterona oscilan con una frecuencia fisiológica, casi horaria; no hace falta una entrada rítmica del hipotálamo para los pulsos. Un modelo hipófisis--suprarrenal con realimentación retardada produce estas oscilaciones. & \cite{Walker2012glc, Walker2010} & hipótesis \\
10 & El destinatario lee la frecuencia de las porciones, no el nivel & Monos sin liberación propia de hormona liberadora de gonadotropinas: la infusión constante no restablece la liberación de hormonas hipofisarias; las porciones una vez por hora sí; volver a la infusión constante suprime de nuevo la respuesta. La fatiga del receptor como filtro pasa altos es una interpretación. & \cite{Belchetz1978} & medido \\
11 & Distintas frecuencias de porciones actúan de forma diferente sobre distintas células de una misma glándula & Los mismos monos: aumentar las porciones a $2$--$5$ por hora reduce ambas hormonas hipofisarias; espaciarlas a una cada $3$~horas reduce una (LH) y aumenta la otra (FSH), de modo que su cociente lo fija la frecuencia. & \cite{Wildt1981gnrh} & medido \\
12 & El ritmo diario lo genera una realimentación negativa intracelular con un retardo de varias horas & Revisión: las proteínas de los genes reloj suprimen con retardo la transcripción de sus propios genes; el lazo de transcripción y traducción da un periodo de alrededor de un día. & \cite{TakahashiJS2017} & medido \\
13 & El generador principal es un grupo de decenas de miles de neuronas que sincroniza el resto del cuerpo & Revisión: el núcleo supraquiasmático es el generador diario principal; sus neuronas aisladas en cultivo son generadores autónomos, y la red las sincroniza; en el ratón, unas $10^4$ neuronas en cada lado. & \cite{Welsh2010scn} & medido \\
14 & El periodo propio en el ser humano es de $24.18$~horas & Personas con luz tenue y un horario desacoplado del día: el periodo de los ritmos de melatonina, temperatura y cortisol es en promedio de $24.18$~horas en jóvenes y mayores, con poca dispersión. & \cite{Czeisler1999} & medido \\
15 & La luz ajusta el generador como un lazo de enganche de fase~\eqref{eq:pll}; hace falta un desplazamiento de unos $11$~min al día & Personas, un solo pulso de luz brillante de $6.7$~horas a distintas horas del día: curva de respuesta de fase de tipo~1 con amplitud de $5$~horas: retraso antes del mínimo de la temperatura corporal, adelanto después. La ecuación~\eqref{eq:pll} es el modelo estándar; $11$~min $=0.18$~horas es aritmética. & \cite{Khalsa2003prc} & medido \\
16 & Sin luz no hay enganche, y el ritmo deriva hacia su periodo propio & Personas totalmente ciegas, sin percepción de luz, que viven en sociedad normal: en cerca de la mitad, los ritmos de melatonina y cortisol siguen su propio periodo, que no coincide con el día. & \cite{Sack1992blind} & medido \\
\bottomrule
\end{longtable}}

\section{Capítulo~\ref{sec:debugging}. Depuración de la máquina}

Qué oye el electrodo, la baja dimensión de la actividad, el funcionamiento de las interfaces con matrices rígidas, el aprendizaje conjunto del cerebro y el decodificador y los puntos visuales producidos por estimulación están medidos en experimentos revisados por pares; las estimaciones de los flujos de datos son aritmética del capítulo; sobre el dispositivo de Neuralink implantado en personas, hasta donde pudimos comprobar, los datos los ha publicado sobre todo la propia empresa.

{\footnotesize
\setlength{\tabcolsep}{3pt}\renewcommand{\arraystretch}{1.15}
\begin{longtable}{>{\raggedright}p{0.5cm}>{\raggedright}p{4.0cm}>{\raggedright}p{5.6cm}>{\raggedright}p{2.3cm}>{\raggedright\arraybackslash}p{1.7cm}}
\toprule
N.º & Afirmación & Experimento y qué se midió & Fuente & Estado \\
\midrule\endfirsthead
\toprule
N.º & Afirmación & Experimento y qué se midió & Fuente & Estado \\
\midrule\endhead
1 & El electrodo oye las espigas de las neuronas en un radio del orden de un centenar de micrómetros, normalmente de una a unas pocas; se distinguen por su forma & Rata, área CA1, registro intracelular y extracelular simultáneo: la forma de la espiga extracelular reproduce la anchura y la amplitud de la intracelular. Estimación: un tetrodo podría distinguir $60$--$100$ neuronas, pero en la práctica se aíslan unas seis. & \cite{Henze2000extra} & medido \\
2 & El cerebro tiene $8.6\cdot10^{10}$ neuronas; la sonda oye más o menos una cienmillonésima & Recuento de núcleos en una suspensión de tejido disuelto: $86\cdot10^9$ neuronas en el cerebro humano. La fracción es aritmética del capítulo. & \cite{Azevedo2009} & medido \\
3 & La actividad de cientos de neuronas de la corteza motora cabe en una o dos decenas de variables comunes & Revisión de registros en la corteza motora de monos: la actividad de la población se describe con pocas variables latentes comunes a muchas neuronas. & \cite{Gallego2017} & medido \\
4 & El ruido de las vecinas está correlacionado, y un centenar de neuronas lleva casi lo mismo que mil (proposición~\ref{prop:pooling}) & Corteza visual de mono: coeficiente de correlación de ruido entre neuronas vecinas de alrededor de $0.1$, y la ganancia por promedio se satura con un centenar de neuronas. Teoría de las correlaciones que limitan la información. & \cite{Zohary1994, MorenoBote2014} & medido \\
5 & Flujo bruto de $2\cdot10^8$~bit/s frente a $8\cdot10^4$~bit/s en espigas~\eqref{eq:probedata} & Aritmética del capítulo basada en la capacidad del flujo de espigas~\eqref{eq:spikeentropy}. Los parámetros de digitalización son reales: el chip de Neuralink trabaja a $19.3$~kHz y $10$~bits por canal. & deducción del capítulo; \cite{Rieke1997, Musk2019} & teoría \\
6 & Primer sistema de Neuralink: los chips amplifican y digitalizan, el flujo bruto va por cable y las espigas las detecta un programa externo; la detección en el chip, la carcasa cerrada, la radio y la alimentación inalámbrica corresponden al dispositivo para personas, según la empresa & Artículo de la empresa: todo el flujo bruto por un cable USB-C; las espigas las detecta en tiempo real un programa fuera del chip; la compresión a bordo, la alimentación y la transmisión inalámbricas figuran como plan para los dispositivos clínicos. Para el dispositivo implantado en personas, solo comunicados de la empresa. Que sea necesario detectar las espigas en el chip es una conclusión del capítulo a partir de~\eqref{eq:probedata}. & \cite{Musk2019}; informes de Neuralink, sin artículo revisado por pares & medido (artículo de la empresa); en personas, informe de la empresa \\
7 & Hasta $3072$ electrodos en $96$ hilos, $32$ por hilo; los hilos los inserta un robot, esquivando los vasos & Artículo de la empresa: $3072$ electrodos en $96$ hilos de $32$; el robot inserta $6$ hilos ($192$ electrodos) por minuto; hasta el $70\%$ de los electrodos oyen espigas en ratas con la matriz implantada a largo plazo. & \cite{Musk2019} & medido (artículo de la empresa) \\
8 & En personas, unos mil canales en $64$ hilos; parte de los hilos se desplazó; cursor del orden de $10$~bit/s & Solo comunicados de la empresa. Una búsqueda en PubMed no encontró ningún artículo revisado por pares con resultados en personas; esto concuerda con la salvedad del texto del capítulo. & informes de Neuralink; sin artículo revisado por pares & medido (informe de la empresa) \\
9 & Una interfaz con una matriz de un centenar de electrodos controla un cursor & Persona con parálisis, $96$ electrodos en la corteza motora primaria: la intención de mover el brazo modifica las espigas tres años después de la lesión; el decodificador da un “cursor neuronal” (correo, televisor), control de una prótesis de mano y de un brazo robótico. & \cite{Hochberg2006} & medido \\
10 & Escritura con la “mano mental”: unos $90$ caracteres por minuto, menos de $8$~bit/s & Persona con parálisis de la mano, intentos de escribir a mano, decodificador recurrente: $90$ caracteres por minuto con una precisión del $94.1\%$ en tiempo real, más del $99\%$ tras la autocorrección. La estimación en bits es aritmética del capítulo. & \cite{Willett2021} & medido \\
11 & Los intentos de hablar se convierten en texto a unas $60$ palabras por minuto & Persona que perdió el habla inteligible, matrices en la corteza: $62$ palabras por minuto; $9.1\%$ de errores de palabra con un vocabulario de $50$ palabras y $23.8\%$ con uno de $125\,000$. & \cite{Willett2023} & medido \\
12 & La interfaz extrae una pequeña fracción de la capacidad: una neurona, decenas de bits por segundo; un centenar, miles (proposición~\ref{prop:probe}) & El límite superior~\eqref{eq:spikeentropy} es teoría; la comparación con las filas~8 y~10 es una conclusión del capítulo. Que el cuello de botella sea la baja dimensión y el desconocimiento del código, y no el cable, no se ha comprobado con un experimento directo. & \cite{Rieke1997} & teoría \\
13 & El cerebro y el decodificador aprenden el uno del otro & Monos controlan un cursor; el decodificador se ajusta en lazo cerrado mientras las neuronas se reaprenden a la vez: la precisión aumenta, la habilidad se conserva y le afectan menos los movimientos del propio brazo. El consejo de separar las velocidades de aprendizaje es una regla de ingeniería, sin experimento propio en el capítulo. & \cite{Orsborn2014coad} & medido \\
14 & La corriente a través de un electrodo excita las neuronas y los cables circundantes sin distinguir tipo & Corteza de ratón y de gato, registro de calcio con dos fotones: incluso una corriente débil excita neuronas escasas, dispersas hasta a milímetros de distancia, mediante la excitación directa de axones en un volumen de decenas de micrómetros de diámetro. Es decir, la excitación no es selectiva, pero tampoco continua. & \cite{Histed2009stim} & medido \\
15 & La estimulación de la corteza visual enciende puntos; hasta $16$ puntos a la vez permiten reconocer algunas letras y los bordes de objetos & Persona ciega, $96$ electrodos en la corteza visual durante $6$ meses: umbral de un electrodo, $66.8\pm36.5$~\textmu A; la estimulación simultánea de grupos (hasta $16$ electrodos, el límite del estimulador) reduce el umbral y da patrones distinguibles; se reconocen algunas letras y bordes de objetos. & \cite{Fernandez2021} & medido \\
16 & La segunda línea de Neuralink es un dispositivo que introduce señales en la corteza visual & Solo comunicados de la empresa sobre su desarrollo; no se encontraron datos revisados por pares sobre su funcionamiento. & informes de Neuralink & hipótesis \\
17 & Las concentraciones de los neurotransmisores de difusión pueden medirse con un sensor químico en el tejido & Rodajas de cerebro de rata, voltamperometría cíclica rápida con fibra de carbono: se midieron la liberación y la recaptación de serotonina; la sustancia sale fuera de la sinapsis. & \cite{Bunin1998} & medido \\
\bottomrule
\end{longtable}}

\section{Capítulo~\ref{sec:eetypes}. Neuronas como instrumentos}

Los modos de trabajo de las células individuales están medidos directamente, con corriente a través de un electrodo y registro de la tensión; la matemática del integrador y la división en clases son teoría clásica, y que el cerebro use la reconfiguración de modos para calcular sigue siendo una hipótesis.

{\footnotesize
\setlength{\tabcolsep}{3pt}\renewcommand{\arraystretch}{1.15}
\begin{longtable}{>{\raggedright}p{0.5cm}>{\raggedright}p{4.0cm}>{\raggedright}p{5.6cm}>{\raggedright}p{2.3cm}>{\raggedright\arraybackslash}p{1.7cm}}
\toprule
N.º & Afirmación & Experimento y qué se midió & Fuente & Estado \\
\midrule\endfirsthead
\toprule
N.º & Afirmación & Experimento y qué se midió & Fuente & Estado \\
\midrule\endhead
1 & Resonador: una corriente lenta que se opone al cambio de tensión convierte la neurona en un filtro pasa banda con máximo entre unos pocos hercios y decenas de hercios (sección~\ref{sec:resonator}) & En rodajas, se inyectó en neuronas \nt{GLUT} una corriente sinusoidal de frecuencia creciente y se midió la impedancia $|Z(f)|$: máximo en $2$--$5$~Hz a $33^\circ$C (unos $7$~Hz a $38^\circ$C); la resonancia la crean corrientes lentas de potasio y catiónica y la amplifica una corriente persistente de sodio. Una revisión muestra el mismo recurso en muchas clases de células & \cite{HuVervaekeStorm2002,HutcheonYarom2000} & medido \\
2 & La corriente lenta se comporta como una inductancia; junto con la capacidad de la membrana forma un circuito resonante & Se midió la respuesta subumbral de un axón a la corriente y se describió con ecuaciones de conductancia linealizadas; el circuito equivalente contiene una inductancia “fenomenológica” cuyo valor depende de la tensión & \cite{Mauro1970} & teoría \\
3 & Constante de tiempo de un grupo con realimentación $\tau_{\text{ef}}=\tau/(1-w)$~\eqref{eq:perfectint}; con $\tau=0.1$~s y $w=0.995$, $20$~s & Se deduce en el capítulo de la ecuación lineal del grupo; como mecanismo para mantener la posición del ojo se propuso en un trabajo teórico & deducción del capítulo; \cite{Seung1996} & teoría \\
4 & El integrador de posición del ojo mantiene la entrada mediante la red, no mediante una propiedad de la célula aislada & Registro intracelular en peces, en el núcleo que mantiene la posición del ojo: tras un giro rápido, la frecuencia de la neurona cambia en escalón y se mantiene; una corriente breve en una sola célula solo causa un desplazamiento pasajero, y los escalones de tensión persisten también por debajo del umbral: los sostienen las entradas sinápticas & \cite{Aksay2001} & medido \\
5 & El integrador sin fugas necesita un reajuste constante mediante aprendizaje; desajustado, olvida o se satura & Se entrenó a peces con una realimentación visual proporcional a $\pm$ la posición del ojo: el mantenimiento se volvía inestable o con fugas, la constante de tiempo de las neuronas caía más de un orden de magnitud, hasta $<1$~s; la realimentación visual normal restablecía poco a poco la estabilidad & \cite{Major2004} & medido \\
6 & Neurona biestable: una entrada breve activa una meseta duradera y la inhibición la apaga; la propiedad la activa la serotonina (sección~\ref{sec:bistable}) & Registro intracelular de neuronas \nt{ACET} que controlan los músculos, en el gato: una excitación sináptica breve lleva la neurona a una actividad duradera y una inhibición breve la reinicia; en el gato espinal, el estado aparece tras administrar un precursor de la serotonina & \cite{Hounsgaard1988} & medido \\
7 & Dos clases: integradores (la frecuencia crece desde cero) y resonadores (en el umbral, de inmediato una frecuencia finita) & Las clases se deducen de la teoría de bifurcaciones. Experimento: en rodajas de corteza, las neuronas \nt{GLUT} empiezan a disparar con una frecuencia tan baja como se quiera (tipo~1), y las neuronas \nt{GABA} rápidas, con un salto a una frecuencia finita (tipo~2) & \cite{Izhikevich2007,Tateno2004} & medido \\
8 & Una misma neurona es integrador o detector de coincidencias: la inhibición acorta la ventana de suma & En rodajas de hipocampo, la ventana en que las entradas excitadoras se suman hasta el umbral es de menos de $2$~ms; la acorta la inhibición directa que llega justo después de la excitación; en las dendritas, donde la inhibición es más débil, la ventana es más ancha & \cite{Pouille2001} & medido \\
9 & Línea de retardo hecha de axones (tabla~\ref{tab:eetypes}) & En la lechuza común se midieron los retardos en los axones que llegan a los detectores de coincidencias: distintas longitudes de axón dan distintos retardos, y los detectores están sintonizados a la diferencia de tiempos de llegada del sonido & \cite{Carr1990} & medido \\
10 & Marcapasos en la vigilia y célula silenciosa en el sueño & Las neuronas \nt{SERO} disparan de forma casi regular durante el día, se frenan en el sueño lento y casi callan en el sueño REM (más detalles en el capítulo~\ref{sec:sleep}) & \cite{JacobsAzmitia1992,McGintyHarper1976} & medido \\
11 & El instrumento lo fijan los canales iónicos y los canales de difusión que los ajustan (proposición~\ref{prop:eetypes}) & Demostrado en redes pequeñas: las sustancias moduladoras cambian las propiedades propias de las neuronas y la fuerza de las sinapsis, de modo que un mismo circuito produce ritmos distintos & \cite{Marder2012} & medido \\
12 & El cerebro usa la reconfiguración de modos para calcular (proposición~\ref{prop:eetypes}) & No hay ningún experimento directo en el que la reconfiguración del modo de una célula sea necesaria para resolver una tarea; solo hay experimentos en los que los moduladores cambian la salida de un circuito & \cite{Marder2012} & hipótesis \\
\bottomrule
\end{longtable}}

\section{Capítulo~\ref{sec:dendrites}. Número de dendritas}

La estructura de las clases y el número de entradas están medidos anatómicamente y con registros eléctricos; la estimación~\eqref{eq:dendriteload} es física elemental del condensador, y la explicación computacional del número de entradas se apoya en teoría y modelos.

{\footnotesize
\setlength{\tabcolsep}{3pt}\renewcommand{\arraystretch}{1.15}
\begin{longtable}{>{\raggedright}p{0.5cm}>{\raggedright}p{4.0cm}>{\raggedright}p{5.6cm}>{\raggedright}p{2.3cm}>{\raggedright\arraybackslash}p{1.7cm}}
\toprule
N.º & Afirmación & Experimento y qué se midió & Fuente & Estado \\
\midrule\endfirsthead
\toprule
N.º & Afirmación & Experimento y qué se midió & Fuente & Estado \\
\midrule\endhead
1 & Capacidad específica de la membrana $c_m\approx1$~\textmu F/cm$^2$ & Se extrajo membrana del cuerpo de la neurona con una pipeta, se aplicó un escalón de voltaje y, por la caída de la corriente capacitiva, se halló la capacidad total, que luego se dividió por el área: $0.9$~\textmu F/cm$^2$ en tres clases de neuronas & \cite{Gentet2000} & medido \\
2 & Tiempo de carga $t\approx c_mA\Delta V/I$~\eqref{eq:dendriteload}: una neurona grande necesita decenas de entradas simultáneas; a una pequeña le basta una sinapsis grande & Estimación según la ley de carga del condensador; los números ($20$ y $300$~pF, $0.1$ y varios nA) son órdenes de magnitud & deducción del capítulo & teoría \\
3 & Una sola sinapsis enorme da una corriente de varios nanoamperios y lleva de inmediato a una neurona pequeña al umbral & Registro simultáneo de la terminal y de la neurona destinataria en una rodaja: corriente de una sinapsis de $2$--$13$~nA; cada espiga de entrada provoca una respuesta supraumbral, con un retardo de unos $0.4$~ms a $36^\circ$C; la destinataria emite \nt{GLYC} & \cite{Borst1995,BorstSoria2012} & medido \\
4 & Cero dendritas: en las neuronas sensoriales del tacto y del dolor, la espiga va del sensor a la médula espinal sin pasar por el cuerpo celular & La estructura (un axón que se divide en dos junto al cuerpo) está establecida anatómicamente. Que la conducción no requiera que el cuerpo sea excitable se mostró con un modelo validado de esta neurona: sin cuerpo excitable, la espiga atraviesa la bifurcación con la misma fiabilidad & \cite{AmirDevor2003} & teoría \\
5 & Una dendrita: la neurona olfativa lleva un solo tipo de receptor y transmite una sola línea de entrada & Revisión de experimentos con genes de receptores: cada neurona olfativa expresa un único gen de receptor de una gran familia & \cite{Mombaerts2004} & medido \\
6 & Dos dendritas: en los detectores de la dirección del sonido, las entradas de los oídos izquierdo y derecho se posan en dendritas distintas & Anatomía y registros en mamíferos, reunidos en una revisión: las entradas de los dos oídos se reparten en dos dendritas opuestas & \cite{Grothe2010} & medido \\
7 & Repartir las entradas en dos dendritas hace el AND más estricto & Mostrado en un modelo: la saturación~\eqref{eq:shunt} en una dendrita impide que las entradas de un solo lado lleguen al umbral, y la detección de coincidencias mejora. No hay ningún experimento directo en el que se cambie la distribución de las entradas & \cite{AgmonSnir1998} & teoría \\
8 & Unos $7\cdot10^{10}$ pequeñas neuronas \nt{GLUT} del circuito de aprendizaje motor, con $\sim4$ entradas cada una & Recuento de células por fraccionamiento isotrópico: en esta parte del cerebro humano hay unos $69$ mil millones de neuronas. Registro de una de estas células en una rata viva: ante un toque llegan ráfagas de corrientes discretas desde pocas entradas, y la célula dispara cuando se suman & \cite{Azevedo2009,Chadderton2004} & medido \\
9 & Un elemento de umbral con $n$ entradas separa como máximo $P_{\max}\approx2n$ patrones aleatorios~\eqref{eq:cover}; para la neurona de salida del circuito de aprendizaje motor, la cota es $\sim2\cdot10^5$, y la estimación realista, $\sim4\cdot10^4$ asociaciones & La cota es un resultado clásico de la geometría combinatoria. La estimación realista es la teoría de capacidad con pesos solo positivos y margen frente al ruido, aplicada al número medido de entradas de esta neurona & \cite{Cover1965,Brunel2004} & teoría \\
10 & Los mezcladores con pocas entradas sirven para expandir la entrada; “cuatro” está cerca del óptimo & Teoría: la dimensión de la representación que da la capa de mezcladores es máxima aproximadamente con el número de entradas que se observa anatómicamente; la hipótesis original de la expansión está en trabajos clásicos & \cite{Marr1969,Albus1971,LitwinKumar2017} & hipótesis \\
11 & Árboles grandes: $10^4$--$10^5$ entradas & Recuento de sinapsis en micrografías electrónicas: $\approx1.7\cdot10^5$ sinapsis en una sola neurona \nt{GABA} de salida del circuito de aprendizaje motor en la rata; unas $30\,000$ entradas excitadoras y $1\,700$ inhibidoras en una neurona \nt{GLUT} del hipocampo & \cite{NapperHarvey1988,Megias2001} & medido \\
12 & Las ramas de un árbol grande son subcircuitos independientes con sus propios umbrales; la neurona es una pequeña red de varias capas & Estimulación puntual en dos lugares de dendritas finas: las entradas en una misma rama se suman de forma sigmoidea, y en ramas distintas, de forma lineal. En las dendritas de neuronas de la corteza humana se hallaron espigas de calcio graduadas que permiten a una sola célula resolver una tarea linealmente no separable. Modelo: para reproducir una neurona hace falta una red profunda & \cite{Polsky2004,Gidon2020,Beniaguev2021} & medido \\
13 & Dendritas-salida: cada rama de una neurona \nt{GABA} de la retina es un detector de dirección independiente & Registro de calcio con dos fotones en ramas individuales: la señal de la rama es selectiva al movimiento del cuerpo celular hacia la punta, pero el voltaje del cuerpo no & \cite{Euler2002} & medido \\
14 & El número de entradas sigue a la tarea (proposición~\ref{prop:dendrites}) & La estructura de las clases está medida (filas 3--13); que el número de entradas se haya elegido por velocidad o por clasificación solo lo muestran la teoría y los modelos para clases concretas & \cite{LitwinKumar2017,AgmonSnir1998} & hipótesis \\
\bottomrule
\end{longtable}}

\section{Capítulo~\ref{sec:sleep}. El sueño}

Los modos del sueño, las repeticiones y la reducción de las sinapsis están medidos con registros de neuronas, microdiálisis y microscopía electrónica; el horario del sueño y el vaivén lento se describen con los modelos del libro, y la función del sueño es sobre todo materia de hipótesis.

{\footnotesize
\setlength{\tabcolsep}{3pt}\renewcommand{\arraystretch}{1.15}
\begin{longtable}{>{\raggedright}p{0.5cm}>{\raggedright}p{4.0cm}>{\raggedright}p{5.6cm}>{\raggedright}p{2.3cm}>{\raggedright\arraybackslash}p{1.7cm}}
\toprule
N.º & Afirmación & Experimento y qué se midió & Fuente & Estado \\
\midrule\endfirsthead
\toprule
N.º & Afirmación & Experimento y qué se midió & Fuente & Estado \\
\midrule\endhead
1 & Durante el sueño las neuronas de la corteza no callan; cambia el modo de trabajo & Registros largos de neuronas corticales en ratas: en el sueño lento la frecuencia sigue siendo distinta de cero, y los periodos de actividad alternan con periodos de silencio; tras una vigilia larga la frecuencia es mayor en todos los estados, y tras el sueño, menor & \cite{Vyazovskiy2009} & medido \\
2 & \nt{ACET} es alto de día y en el sueño REM, y bajo en el sueño lento (tabla~\ref{tab:sleepmodes}) & Microdiálisis en la corteza de gatos en libre movimiento: en la vigilia la liberación de acetilcolina es un $100$--$175\%$ mayor que en el sueño lento, y en el sueño REM está más o menos al nivel de la vigilia tranquila & \cite{Marrosu1995} & medido \\
3 & \nt{NORA} y \nt{SERO} se apagan en el sueño lento y casi callan en el sueño REM & Registro de neuronas \nt{NORA} individuales en ratas y de neuronas \nt{SERO} en gatos por fases del sueño: la frecuencia es máxima en la vigilia, menor en el sueño lento y casi nula en el sueño REM & \cite{AstonJonesBloom1981,McGintyHarper1976} & medido \\
4 & En el sueño REM los músculos están desconectados por inhibición a través de \nt{GABA} y \nt{GLYC} & En ratas se midió la actividad de las neuronas \nt{ACET} del músculo masetero y se aplicaron bloqueantes directamente sobre ellas: la parálisis solo desaparece si se bloquean a la vez los receptores \nt{GABA} de ambos tipos y los receptores \nt{GLYC} & \cite{BrooksPeever2012} & medido \\
5 & La presión de sueño $S$ es un condensador con dos umbrales~\eqref{eq:sleeppressure}, $\tau_r=18.2$~h, $\tau_d=4.2$~h & Modelo de dos procesos; las constantes de tiempo se obtienen de cómo la potencia de las ondas lentas del EEG crece con la vigilia y cae durante el sueño, y la forma de los umbrales, del momento del despertar espontáneo & \cite{Borbely1982,Daan1984} & teoría \\
6 & La máquina llega por sí sola a $16.1$~h de vigilia y $8.0$~h de sueño (fig.~\ref{fig:twoprocess}) & Cálculo según~\eqref{eq:sleeppressure} con umbrales oscilantes, script \texttt{models/sleep\_models.py}: $16.1$~h de vigilia y $7.9$--$8.0$~h de sueño & --- & modelo \\
7 & Tras $40$~h sin dormir, $S=0.90$, pero el sueño dura solo $8.8$~h: la deuda se devuelve con profundidad, no con duración & Modelo: \texttt{models/sleep\_models.py} da $S=0.90$ y $8.8$~h. Experimento: tras $40.5$~h sin dormir, en ocho personas la potencia de las ondas lentas del EEG en la primera noche de recuperación es significativamente mayor de lo habitual & \cite{Borbely1981} & modelo \\
8 & Físicamente, $S$ es la adenosina & Microdiálisis en gatos: la adenosina extracelular en una de las regiones que controlan la vigilia aumenta durante la vigilia, sigue aumentando con la privación de sueño y cae durante el sueño. Que $S$ sea solo la adenosina no está demostrado & \cite{PorkkaHeiskanen1997,Dunwiddie2001} & hipótesis \\
9 & En el sueño lento la corteza oscila: actividad durante unos cientos de ms y silencio, menos de una vez por segundo ($<1$~Hz; bajo anestesia, sobre todo $0.3$--$0.4$~Hz) & Registro intracelular en gatos anestesiados: despolarización de $0.8$--$1.5$~s y una larga hiperpolarización, ritmo $<1$~Hz, sobre todo $0.3$--$0.4$~Hz; la misma alternancia de actividad y silencio se ve en el sueño natural (fila~1) & \cite{Steriade1993,Vyazovskiy2009} & medido \\
10 & El vaivén es un grupo con realimentación y fatiga~\eqref{eq:updown}; con $b=5$, periodo de $1.1$~s (valor del modelo) y actividad alrededor del $35\%$ del tiempo; con $b=1$, actividad uniforme (fig.~\ref{fig:updown}) & Cálculo, script \texttt{models/sleep\_models.py}: periodo de $1.11$~s, fracción de tiempo activo $0.35$ (promedio sobre un número entero de periodos); con $b=1$ la fracción activa es $1.0$. El periodo del modelo es más corto que el medido (fila~9) & --- & modelo \\
11 & \nt{ACET} apaga el vaivén y pasa la corteza a un trabajo uniforme & La estimulación de un núcleo de neuronas \nt{ACET} en el gato bloqueó el ritmo lento en el $79\%$ de las neuronas corticales y las pasó a trabajo uniforme, sobre todo porque desaparecieron las hiperpolarizaciones largas. La acetilcolina suprime las corrientes lentas de potasio que producen la fatiga & \cite{SteriadeAmzica1993,McCormick1992} & medido \\
12 & Las ráfagas de ritmo de $7$--$15$~Hz las genera el lazo \nt{GLUT}--\nt{GABA} entre la corteza y un nodo de relevo & En rodajas del nodo de relevo visual del hurón, las ráfagas las producen las ráfagas de rebote de las células de relevo ante la inhibición de un núcleo \nt{GABA} vecino, al que ellas mismas excitan & \cite{vonKrosigk1993,SteriadeMcCormickSejnowski1993} & medido \\
13 & Las repeticiones del día van en avance rápido: en el hipocampo, unas $20$ veces más rápido; en la corteza, $6$--$7$ veces & En el sueño lento, las secuencias repetidas de neuronas del hipocampo van comprimidas en el tiempo. Tras correr por una pista, las secuencias de neuronas de lugar se repitieron en el mismo orden en ráfagas breves de $\sim100$~ms, comprimidas unas $20$ veces; en la corteza, la compresión es de $6$--$7$ veces & \cite{Nadasdy1999,LeeWilson2002,Euston2007} & medido \\
14 & Reforzar la coordinación de las repeticiones con la corteza mejora la memoria (relación causal) & En ratas, tras el aprendizaje, una estimulación ligada a las repeticiones reforzó su acoplamiento con las ondas lentas y las ráfagas de ritmo de la corteza: al día siguiente la memoria era alta, y en los controles, al nivel del azar & \cite{Maingret2016} & medido \\
15 & La función de las repeticiones es el aprendizaje intercalado de la memoria lenta & Teoría de los dos sistemas de aprendizaje y cálculos con redes artificiales: el aprendizaje secuencial estropea lo antiguo, el intercalado no. No hay ningún experimento directo en el cerebro de que las repeticiones sirvan precisamente para eso & \cite{McClelland1995} & hipótesis \\
16 & Tras el sueño, las sinapsis se reducen en proporción a su tamaño; las grandes casi no cambian & Microscopía electrónica tridimensional de $6920$ sinapsis corticales de ratón: el área de contacto tras el sueño es $\sim18\%$ menor que tras la vigilia; la reducción es proporcional al tamaño, y las sinapsis grandes y estables no cambian & \cite{deVivo2017} & medido \\
17 & La tarea principal del sueño es la renormalización de los pesos & Hipótesis de la homeostasis sináptica; las filas 1 y 16 concuerdan con ella, pero no demuestran que sea la función principal & \cite{TononiCirelli2014} & hipótesis \\
18 & El sueño REM es una prueba en banco del modelo del mundo & El modo (tabla~\ref{tab:sleepmodes}) está medido; que la red ejecute el modelo del mundo es una suposición teórica, sin experimento & \cite{Hobson2009} & hipótesis \\
19 & Limpieza del cerebro durante el sueño: cuestión abierta & Un experimento: durante el sueño, el espacio intercelular es un $60\%$ más ancho y la limpieza es más rápida. Otro, con un método de medida distinto: durante el sueño y bajo anestesia, la limpieza es notablemente más lenta. Los resultados se contradicen & \cite{Xie2013,Miao2024} & hipótesis \\
20 & Sueño local: zonas de la corteza de un animal despierto caen brevemente en silencio & En ratas tras una vigilia larga, las neuronas de una zona aislada de la corteza callan brevemente, como en el sueño lento, con una onda lenta en el EEG local; en esos momentos la rata se equivoca más en la tarea & \cite{Vyazovskiy2011} & medido \\
\bottomrule
\end{longtable}}

\section{Capítulo~\ref{sec:livingmachine}. Una máquina de neuronas vivas}

El comportamiento de los cultivos sobre matrices de electrodos está medido en experimentos directos con registro y estimulación; el mecanismo de las ráfagas se apoya en el modelo de agotamiento sináptico y en experimentos con microcultivos, y las afirmaciones sobre la dopamina en la placa, en experimentos aislados, algunos de los cuales los propios autores consideran solo una demostración.

{\footnotesize
\setlength{\tabcolsep}{3pt}\renewcommand{\arraystretch}{1.15}
\begin{longtable}{>{\raggedright}p{0.5cm}>{\raggedright}p{4.0cm}>{\raggedright}p{5.6cm}>{\raggedright}p{2.3cm}>{\raggedright\arraybackslash}p{1.7cm}}
\toprule
N.º & Afirmación & Experimento y qué se midió & Fuente & Estado \\
\midrule\endfirsthead
\toprule
N.º & Afirmación & Experimento y qué se midió & Fuente & Estado \\
\midrule\endhead
1 & Cultivo sobre chip: sobre todo \nt{GLUT}, con una fracción menor de \nt{GABA} que depende de la preparación; en cultivos de hipocampo, alrededor del $6\%$ & Una red de miles de neuronas sobre una matriz de electrodos es un experimento estándar (fila~3). La fracción de neuronas \nt{GABA} en cultivos de hipocampo se contó por tinción de GABA: alrededor del $6\%$ de las neuronas. Para cultivos de corteza no damos una cifra verificada & \cite{Benson1994} & medido \\
2 & Organoides: tejido nervioso tridimensional de alrededor de un milímetro, a partir de células madre humanas & A partir de células madre se cultivaron organoides tridimensionales con varias regiones cerebrales, entre ellas corteza con capas de progenitores y neuronas maduras & \cite{Lancaster2013} & medido \\
3 & Sin entrada, el cultivo se enciende en ráfagas con intervalos de un segundo a minutos, según el cultivo & Se registraron durante cinco semanas $58$ cultivos de corteza con densidades de $3\,000$ a $50\,000$ neuronas ($963$ registros de media hora): tras dos semanas, en casi todos los cultivos predominan las ráfagas de toda la red; los intervalos entre ráfagas van de $1$ a $300$~s y dependen mucho de la preparación & \cite{Wagenaar2006} & medido \\
4 & La ráfaga termina por agotamiento del neurotransmisor, con intervalo $T_{\text{ráfaga}}\approx\tau_{\text{rec}}\ln\frac{1}{1-x^*}$~\eqref{eq:burstinterval}; $2$--$4$~s es la estimación del modelo con $\tau_{\text{rec}}=0.8$~s, la recuperación medida es más lenta & La fórmula se deduce en el capítulo. Modelo: una red con sinapsis que se agotan genera por sí sola ráfagas de toda la red. Experimento en microcultivos de $4$--$30$ neuronas: la duración de la ráfaga depende del tiempo transcurrido desde la anterior, y el agotamiento de las vesículas de neurotransmisor suprime las ráfagas; conclusión de los autores: la ráfaga está limitada por la reserva de neurotransmisor. Allí la recuperación de la reserva tarda decenas de segundos, lo que con la misma fórmula da intervalos de decenas de segundos y minutos & deducción en el capítulo; \cite{Tsodyks2000,CohenSegal2011,Tsodyks1997} & teoría \\
5 & “Un maestro que se calla”: la estimulación se apaga cuando aparece la respuesta deseada, y la respuesta se fija & Se estimuló el cultivo a $0.3$--$1$~Hz hasta que aparecía la respuesta elegida $50\pm10$~ms después del estímulo; luego se apagaba la estimulación durante $5$~min; tras varios ciclos la respuesta deseada se provocaba enseguida; se trazaron curvas de aprendizaje & \cite{Shahaf2001} & medido \\
6 & El cultivo juega al tenis; un aumento significativo de las series en la sesión, solo con retroalimentación completa & En detalle, en el capítulo~\ref{sec:dishbrain} y su apéndice: con silencio en lugar del estímulo, el cultivo al final de la sesión también es mejor que sin retroalimentación, pero con un efecto menor; el aprendizaje de una sesión a otra es inestable & \cite{Kagan2022} & medido \\
7 & Que el cultivo aprenda a predecir su entrada es una hipótesis; las palabras grandilocuentes sobre su “sintiencia” están cuestionadas & El principio de energía libre es una propuesta teórica; no hay un experimento directo que lo distinga de explicaciones más sencillas. Treinta investigadores señalaron en un artículo de respuesta que un cambio de comportamiento en el lazo no muestra ni inteligencia ni sensaciones & \cite{Friston2010,Balci2023} & hipótesis \\
8 & El cultivo conduce un cuerpo virtual hacia un objetivo con un patrón de estímulos elegido & Un programa elegía por sí mismo los patrones de estimulación de entrenamiento según el resultado actual; en decenas de minutos la red alcanzaba los estados deseados, y la plasticidad se mantenía $80$~min tras $2$~h de entrenamiento. Los mismos estímulos, aplicados sin relación con el comportamiento, no tenían efecto & \cite{Bakkum2008} & medido \\
9 & Reservorio: solo se entrena la lectura lineal $y=W_{\text{sal}}\mathbf r$~\eqref{eq:reservoir} & Teoría de la computación de reservorio: cualquier sistema no lineal con memoria que se desvanece más una salida lineal entrenada & \cite{Maass2002,JaegerHaas2004} & teoría \\
10 & Un organoide como reservorio distingue hablantes con una exactitud de alrededor del $78\%$ (según los autores); por el error aprende la lectura, y las conexiones del organoide cambian sin maestro & Sonidos del habla de varios hablantes se entregaron al organoide como patrones de corriente en la matriz de electrodos; la lectura entrenada distinguía al hablante con una exactitud de alrededor del $78\%$, según informan los autores. Ellos mismos informan que durante el entrenamiento cambiaba la conectividad funcional del organoide, y lo llaman aprendizaje no supervisado en el propio organoide & \cite{Cai2023} & medido \\
11 & Un millón de neuronas gasta en espigas unos $10^{-4}$~W~\eqref{eq:dishpower} & Estimación: el coste de una espiga, $10^{-10}$~J, del presupuesto energético de la corteza~\eqref{eq:energy}, multiplicado por el número de espigas; no hay medición directa de la potencia del cultivo & deducción en el capítulo; \cite{Attwell2001} & teoría \\
12 & Dopamina con un destello de luz: $1$~mM de dopamina enjaulada, $240$ repeticiones, el número de espigas evocadas crecía & En una plataforma remota con organoides, la dopamina en una jaula fotosensible ($1$~mM) se liberaba con ultravioleta ($365$~nm, $0.8$~s) cuando el estímulo provocaba más espigas; a lo largo de $240$ pasos (unas $5$~h) el número de espigas creció en conjunto. Los autores escriben que con un solo experimento no se puede establecer la relación con la dopamina; no hay controles & \cite{Jordan2024} & hipótesis \\
13 & La dopamina de células vivas cambia el peso de un transistor orgánico de forma duradera y reversible & Células que liberan dopamina se cultivaron sobre un elemento neuromórfico orgánico; la oxidación de la dopamina en el electrodo cambia la conductancia del canal; se mostraron un cambio duradero del peso y su retorno & \cite{Keene2020} & medido \\
14 & Existen organoides con neuronas \nt{DOPA} funcionales; su dopamina no se ha usado como maestro & En organoides de células madre humanas se encontraron neuronas \nt{DOPA} maduras eléctricamente activas y producción de dopamina. No hemos encontrado en la literatura experimentos en los que esta dopamina enseñe a otra red & \cite{Jo2016} & medido \\
15 & Un organoide equilibra un péndulo: los estímulos de enseñanza del programa son mejores que los aleatorios, sin consolidación, a través de sinapsis \nt{GLUT} & Organoides de corteza de ratón en lazo cerrado con la tarea del “péndulo sobre un carro”: en la mayoría de los organoides, las señales elegidas por aprendizaje por refuerzo dieron mejor resultado que las aleatorias o que ninguna; tras $45$~min de descanso la mejora no se conservaba; el bloqueo de los receptores AMPA y NMDA la anulaba & \cite{Robbins2026} & medido \\
16 & Sin registros de control, la red del banco no puede decidir por sí misma qué cuenta como éxito (proposición~\ref{prop:livingmachine}) & En todos los experimentos de las filas 5--15 la señal de enseñanza la fija el experimentador o un programa; no hay ningún experimento en el que el cultivo haya elaborado por sí mismo un criterio de éxito: es una conclusión por ausencia, no un experimento & --- & hipótesis \\
\bottomrule
\end{longtable}}

\section{Capítulo~\ref{sec:dishbrain}. DishBrain}

El diseño del banco y los resultados del juego provienen de un solo trabajo experimental y de su continuación; las fórmulas del ruido y de la deriva hacia los buenos ajustes se deducen en el capítulo y se comprueban con el modelo del libro, y el mecanismo del aprendizaje en la placa no está establecido.

{\footnotesize
\setlength{\tabcolsep}{3pt}\renewcommand{\arraystretch}{1.15}
\begin{longtable}{>{\raggedright}p{0.5cm}>{\raggedright}p{4.0cm}>{\raggedright}p{5.6cm}>{\raggedright}p{2.3cm}>{\raggedright\arraybackslash}p{1.7cm}}
\toprule
N.º & Afirmación & Experimento y qué se midió & Fuente & Estado \\
\midrule\endfirsthead
\toprule
N.º & Afirmación & Experimento y qué se midió & Fuente & Estado \\
\midrule\endhead
1 & Banco: unas $800\,000$ células de corteza de ratón o unas $10^6$ células humanas por chip; $26\,000$ electrodos en $8$~mm$^2$, registro de hasta $1024$, estimulación a través de $8$ & Métodos del trabajo: chip MaxOne, $26\,000$ electrodos de platino en $8$~mm$^2$, registro de $1024$ canales a $20\,000$ muestras por segundo; técnicamente se pueden estimular hasta $32$ electrodos, de forma independiente se logró con $8$; los cultivos de ratón se usaron desde unos $10$~días & \cite{Kagan2022} & medido \\
2 & Lazo con ciclo de $10$~ms: las espigas de las regiones motoras mueven la raqueta (fig.~\ref{fig:dishloop}) & Las espigas se cuentan en ventanas de $10$~ms; el retardo desde la espiga hasta el estímulo de respuesta es de unos $5$~ms y lo fija el búfer del equipo & \cite{Kagan2022} & medido \\
3 & Entrada: código de lugar (cuál de los $8$ electrodos) y de frecuencia, $4$--$40$~pulsos/s & En el experimento principal la frecuencia de estimulación crece linealmente de $4$~Hz, con la pelota junto a la pared lejana, a $40$~Hz junto a la raqueta; la amplitud de los pulsos de la pelota es de $75$~mV; los pulsos son rectangulares bifásicos & \cite{Kagan2022} & medido \\
4 & Salida $\Delta p=c\,(n_\uparrow-n_\downarrow)$~\eqref{eq:dishreadout}, ruido de la cuenta por ventana $1/\sqrt{RT}$ & La regla de la diferencia entre dos regiones está descrita en el trabajo (con pesos de las regiones según su actividad), pero no se da la fórmula exacta. La estimación del ruido es consecuencia del conteo de Poisson y se deduce en el capítulo & \cite{Kagan2022}; deducción en el capítulo & teoría \\
5 & Maestro: tras el acierto, un estímulo predecible de $75$~mV, $100$~pulsos/s, $0.1$~s; tras el fallo, uno impredecible, $150$~mV, $5$~pulsos/s, $4$~s en lugares e instantes aleatorios, y luego $4$~s de pausa & Métodos: tras el acierto, $75$~mV, $100$~Hz, $100$~ms en los $8$ electrodos a la vez; tras el fallo, $150$~mV, $5$~Hz en electrodos aleatorios e instantes aleatorios durante $4$~s, y luego $4$~s sin estimulación. En el cultivo no hay dopamina & \cite{Kagan2022} & medido \\
6 & La red actúa de modo que su entrada se vuelva predecible & El principio de energía libre es una propuesta teórica de la que los autores derivaron el protocolo; el experimento concuerda con él, pero no lo distingue de mecanismos más sencillos (filas~7--8) & \cite{Friston2010,Kagan2022} & hipótesis \\
7 & Si el fracaso sacude la red y el éxito no, el tiempo en un ajuste es $\rho\propto1/P_{\text{fallo}}$~\eqref{eq:dishdensity} (proposición~\ref{prop:dishscramble}) & Se sigue del balance de flujos en una cadena de Markov con empujones simétricos; se deduce en el capítulo. No hay comprobación directa en un cultivo & deducción en el capítulo & teoría \\
8 & El modelo de “revolver tras el fallo” aprende: fracción de aciertos $0.21\to0.38$; con empujón tras cada peloteo, $\approx0.12$; sin empujones, $0.19$ (fig.~\ref{fig:dishmodel}) & Cálculo, script \texttt{models/dishbrain\_model.py}, $40$ cultivos modelo con $2000$ peloteos cada uno: $0.21\to0.38$; $0.13\to0.11$; $0.19\to0.19$ & --- & modelo \\
9 & Las series de pelotas devueltas se alargan solo en cultivos vivos con el lazo completo; los controles no aprenden & $399$ sesiones de $20$~min: $9$ cultivos de ratón ($101$ sesiones), $11$ humanos ($138$), $6$ chips con medio ($80$), $20$ cultivos en reposo ($42$), $3$ simulaciones ($38$). La longitud media del peloteo en los últimos $15$~min es mayor que en los primeros $5$ solo en cultivos vivos; con células que no son neuronas (HEK293T) y con chips con medio no hay efecto & \cite{Kagan2022} & medido \\
10 & Un aumento significativo en la sesión, solo con realimentación completa; con silencio en lugar del estímulo, al final de la sesión se juega mejor que sin realimentación, pero el efecto es menor & $486$ sesiones en $3$ días: “estímulo” ($n=27$), “silencio” ($17$), “sin realimentación” ($15$). El aumento en el tiempo solo es significativo en la condición “estímulo”; al final de la sesión, “silencio” superaba a “sin realimentación” con un efecto menor & \cite{Kagan2022} & medido \\
11 & El efecto es pequeño; si la red aprende de forma duradera es una cuestión abierta & Dentro de la sesión la mejora es sólida, pero, según los propios autores, no se observó un aprendizaje estable entre sesiones a lo largo de varios días & \cite{Kagan2022} & medido \\
12 & Durante el juego la red se acerca al régimen crítico, y cuanto más cerca, mejor el juego & Análisis de avalanchas de actividad en los mismos cultivos: la entrada estructurada acerca la red al régimen crítico, y la calidad del juego se correlaciona con la cercanía a él; pero la criticidad sola, sin información sobre las consecuencias de las acciones, no basta para aprender & \cite{Habibollahi2023} & medido \\
13 & La “sintiencia” del cultivo no está demostrada & Artículo de respuesta de treinta investigadores: un cambio de comportamiento en lazo cerrado no demuestra ni inteligencia ni sensaciones; no hay un experimento que lo muestre & \cite{Balci2023} & hipótesis \\
14 & Cuarto control propuesto: un estímulo predecible tras el fallo, de la misma duración y energía (sección~\ref{sec:dishspec}) & Este experimento no se ha hecho; es una propuesta del libro & --- & hipótesis \\
\bottomrule
\end{longtable}}

